\documentclass{article} % For LaTeX2e
\usepackage{iclr2027_conference,times}

% Optional math commands from https://github.com/goodfeli/dlbook_notation.
%%%%% NEW MATH DEFINITIONS %%%%%

\usepackage{amsmath,amsfonts,bm}

% Mark sections of captions for referring to divisions of figures
\newcommand{\figleft}{{\em (Left)}}
\newcommand{\figcenter}{{\em (Center)}}
\newcommand{\figright}{{\em (Right)}}
\newcommand{\figtop}{{\em (Top)}}
\newcommand{\figbottom}{{\em (Bottom)}}
\newcommand{\captiona}{{\em (a)}}
\newcommand{\captionb}{{\em (b)}}
\newcommand{\captionc}{{\em (c)}}
\newcommand{\captiond}{{\em (d)}}

% Highlight a newly defined term
\newcommand{\newterm}[1]{{\bf #1}}


% Figure reference, lower-case.
\def\figref#1{figure~\ref{#1}}
% Figure reference, capital. For start of sentence
\def\Figref#1{Figure~\ref{#1}}
\def\twofigref#1#2{figures \ref{#1} and \ref{#2}}
\def\quadfigref#1#2#3#4{figures \ref{#1}, \ref{#2}, \ref{#3} and \ref{#4}}
% Section reference, lower-case.
\def\secref#1{section~\ref{#1}}
% Section reference, capital.
\def\Secref#1{Section~\ref{#1}}
% Reference to two sections.
\def\twosecrefs#1#2{sections \ref{#1} and \ref{#2}}
% Reference to three sections.
\def\secrefs#1#2#3{sections \ref{#1}, \ref{#2} and \ref{#3}}
% Reference to an equation, lower-case.
\def\eqref#1{equation~\ref{#1}}
% Reference to an equation, upper case
\def\Eqref#1{Equation~\ref{#1}}
% A raw reference to an equation---avoid using if possible
\def\plaineqref#1{\ref{#1}}
% Reference to a chapter, lower-case.
\def\chapref#1{chapter~\ref{#1}}
% Reference to an equation, upper case.
\def\Chapref#1{Chapter~\ref{#1}}
% Reference to a range of chapters
\def\rangechapref#1#2{chapters\ref{#1}--\ref{#2}}
% Reference to an algorithm, lower-case.
\def\algref#1{algorithm~\ref{#1}}
% Reference to an algorithm, upper case.
\def\Algref#1{Algorithm~\ref{#1}}
\def\twoalgref#1#2{algorithms \ref{#1} and \ref{#2}}
\def\Twoalgref#1#2{Algorithms \ref{#1} and \ref{#2}}
% Reference to a part, lower case
\def\partref#1{part~\ref{#1}}
% Reference to a part, upper case
\def\Partref#1{Part~\ref{#1}}
\def\twopartref#1#2{parts \ref{#1} and \ref{#2}}

\def\ceil#1{\lceil #1 \rceil}
\def\floor#1{\lfloor #1 \rfloor}
\def\1{\bm{1}}
\newcommand{\train}{\mathcal{D}}
\newcommand{\valid}{\mathcal{D_{\mathrm{valid}}}}
\newcommand{\test}{\mathcal{D_{\mathrm{test}}}}

\def\eps{{\epsilon}}


% Random variables
\def\reta{{\textnormal{$\eta$}}}
\def\ra{{\textnormal{a}}}
\def\rb{{\textnormal{b}}}
\def\rc{{\textnormal{c}}}
\def\rd{{\textnormal{d}}}
\def\re{{\textnormal{e}}}
\def\rf{{\textnormal{f}}}
\def\rg{{\textnormal{g}}}
\def\rh{{\textnormal{h}}}
\def\ri{{\textnormal{i}}}
\def\rj{{\textnormal{j}}}
\def\rk{{\textnormal{k}}}
\def\rl{{\textnormal{l}}}
% rm is already a command, just don't name any random variables m
\def\rn{{\textnormal{n}}}
\def\ro{{\textnormal{o}}}
\def\rp{{\textnormal{p}}}
\def\rq{{\textnormal{q}}}
\def\rr{{\textnormal{r}}}
\def\rs{{\textnormal{s}}}
\def\rt{{\textnormal{t}}}
\def\ru{{\textnormal{u}}}
\def\rv{{\textnormal{v}}}
\def\rw{{\textnormal{w}}}
\def\rx{{\textnormal{x}}}
\def\ry{{\textnormal{y}}}
\def\rz{{\textnormal{z}}}

% Random vectors
\def\rvepsilon{{\mathbf{\epsilon}}}
\def\rvtheta{{\mathbf{\theta}}}
\def\rva{{\mathbf{a}}}
\def\rvb{{\mathbf{b}}}
\def\rvc{{\mathbf{c}}}
\def\rvd{{\mathbf{d}}}
\def\rve{{\mathbf{e}}}
\def\rvf{{\mathbf{f}}}
\def\rvg{{\mathbf{g}}}
\def\rvh{{\mathbf{h}}}
\def\rvu{{\mathbf{i}}}
\def\rvj{{\mathbf{j}}}
\def\rvk{{\mathbf{k}}}
\def\rvl{{\mathbf{l}}}
\def\rvm{{\mathbf{m}}}
\def\rvn{{\mathbf{n}}}
\def\rvo{{\mathbf{o}}}
\def\rvp{{\mathbf{p}}}
\def\rvq{{\mathbf{q}}}
\def\rvr{{\mathbf{r}}}
\def\rvs{{\mathbf{s}}}
\def\rvt{{\mathbf{t}}}
\def\rvu{{\mathbf{u}}}
\def\rvv{{\mathbf{v}}}
\def\rvw{{\mathbf{w}}}
\def\rvx{{\mathbf{x}}}
\def\rvy{{\mathbf{y}}}
\def\rvz{{\mathbf{z}}}

% Elements of random vectors
\def\erva{{\textnormal{a}}}
\def\ervb{{\textnormal{b}}}
\def\ervc{{\textnormal{c}}}
\def\ervd{{\textnormal{d}}}
\def\erve{{\textnormal{e}}}
\def\ervf{{\textnormal{f}}}
\def\ervg{{\textnormal{g}}}
\def\ervh{{\textnormal{h}}}
\def\ervi{{\textnormal{i}}}
\def\ervj{{\textnormal{j}}}
\def\ervk{{\textnormal{k}}}
\def\ervl{{\textnormal{l}}}
\def\ervm{{\textnormal{m}}}
\def\ervn{{\textnormal{n}}}
\def\ervo{{\textnormal{o}}}
\def\ervp{{\textnormal{p}}}
\def\ervq{{\textnormal{q}}}
\def\ervr{{\textnormal{r}}}
\def\ervs{{\textnormal{s}}}
\def\ervt{{\textnormal{t}}}
\def\ervu{{\textnormal{u}}}
\def\ervv{{\textnormal{v}}}
\def\ervw{{\textnormal{w}}}
\def\ervx{{\textnormal{x}}}
\def\ervy{{\textnormal{y}}}
\def\ervz{{\textnormal{z}}}

% Random matrices
\def\rmA{{\mathbf{A}}}
\def\rmB{{\mathbf{B}}}
\def\rmC{{\mathbf{C}}}
\def\rmD{{\mathbf{D}}}
\def\rmE{{\mathbf{E}}}
\def\rmF{{\mathbf{F}}}
\def\rmG{{\mathbf{G}}}
\def\rmH{{\mathbf{H}}}
\def\rmI{{\mathbf{I}}}
\def\rmJ{{\mathbf{J}}}
\def\rmK{{\mathbf{K}}}
\def\rmL{{\mathbf{L}}}
\def\rmM{{\mathbf{M}}}
\def\rmN{{\mathbf{N}}}
\def\rmO{{\mathbf{O}}}
\def\rmP{{\mathbf{P}}}
\def\rmQ{{\mathbf{Q}}}
\def\rmR{{\mathbf{R}}}
\def\rmS{{\mathbf{S}}}
\def\rmT{{\mathbf{T}}}
\def\rmU{{\mathbf{U}}}
\def\rmV{{\mathbf{V}}}
\def\rmW{{\mathbf{W}}}
\def\rmX{{\mathbf{X}}}
\def\rmY{{\mathbf{Y}}}
\def\rmZ{{\mathbf{Z}}}

% Elements of random matrices
\def\ermA{{\textnormal{A}}}
\def\ermB{{\textnormal{B}}}
\def\ermC{{\textnormal{C}}}
\def\ermD{{\textnormal{D}}}
\def\ermE{{\textnormal{E}}}
\def\ermF{{\textnormal{F}}}
\def\ermG{{\textnormal{G}}}
\def\ermH{{\textnormal{H}}}
\def\ermI{{\textnormal{I}}}
\def\ermJ{{\textnormal{J}}}
\def\ermK{{\textnormal{K}}}
\def\ermL{{\textnormal{L}}}
\def\ermM{{\textnormal{M}}}
\def\ermN{{\textnormal{N}}}
\def\ermO{{\textnormal{O}}}
\def\ermP{{\textnormal{P}}}
\def\ermQ{{\textnormal{Q}}}
\def\ermR{{\textnormal{R}}}
\def\ermS{{\textnormal{S}}}
\def\ermT{{\textnormal{T}}}
\def\ermU{{\textnormal{U}}}
\def\ermV{{\textnormal{V}}}
\def\ermW{{\textnormal{W}}}
\def\ermX{{\textnormal{X}}}
\def\ermY{{\textnormal{Y}}}
\def\ermZ{{\textnormal{Z}}}

% Vectors
\def\vzero{{\bm{0}}}
\def\vone{{\bm{1}}}
\def\vmu{{\bm{\mu}}}
\def\vtheta{{\bm{\theta}}}
\def\va{{\bm{a}}}
\def\vb{{\bm{b}}}
\def\vc{{\bm{c}}}
\def\vd{{\bm{d}}}
\def\ve{{\bm{e}}}
\def\vf{{\bm{f}}}
\def\vg{{\bm{g}}}
\def\vh{{\bm{h}}}
\def\vi{{\bm{i}}}
\def\vj{{\bm{j}}}
\def\vk{{\bm{k}}}
\def\vl{{\bm{l}}}
\def\vm{{\bm{m}}}
\def\vn{{\bm{n}}}
\def\vo{{\bm{o}}}
\def\vp{{\bm{p}}}
\def\vq{{\bm{q}}}
\def\vr{{\bm{r}}}
\def\vs{{\bm{s}}}
\def\vt{{\bm{t}}}
\def\vu{{\bm{u}}}
\def\vv{{\bm{v}}}
\def\vw{{\bm{w}}}
\def\vx{{\bm{x}}}
\def\vy{{\bm{y}}}
\def\vz{{\bm{z}}}

% Elements of vectors
\def\evalpha{{\alpha}}
\def\evbeta{{\beta}}
\def\evepsilon{{\epsilon}}
\def\evlambda{{\lambda}}
\def\evomega{{\omega}}
\def\evmu{{\mu}}
\def\evpsi{{\psi}}
\def\evsigma{{\sigma}}
\def\evtheta{{\theta}}
\def\eva{{a}}
\def\evb{{b}}
\def\evc{{c}}
\def\evd{{d}}
\def\eve{{e}}
\def\evf{{f}}
\def\evg{{g}}
\def\evh{{h}}
\def\evi{{i}}
\def\evj{{j}}
\def\evk{{k}}
\def\evl{{l}}
\def\evm{{m}}
\def\evn{{n}}
\def\evo{{o}}
\def\evp{{p}}
\def\evq{{q}}
\def\evr{{r}}
\def\evs{{s}}
\def\evt{{t}}
\def\evu{{u}}
\def\evv{{v}}
\def\evw{{w}}
\def\evx{{x}}
\def\evy{{y}}
\def\evz{{z}}

% Matrix
\def\mA{{\bm{A}}}
\def\mB{{\bm{B}}}
\def\mC{{\bm{C}}}
\def\mD{{\bm{D}}}
\def\mE{{\bm{E}}}
\def\mF{{\bm{F}}}
\def\mG{{\bm{G}}}
\def\mH{{\bm{H}}}
\def\mI{{\bm{I}}}
\def\mJ{{\bm{J}}}
\def\mK{{\bm{K}}}
\def\mL{{\bm{L}}}
\def\mM{{\bm{M}}}
\def\mN{{\bm{N}}}
\def\mO{{\bm{O}}}
\def\mP{{\bm{P}}}
\def\mQ{{\bm{Q}}}
\def\mR{{\bm{R}}}
\def\mS{{\bm{S}}}
\def\mT{{\bm{T}}}
\def\mU{{\bm{U}}}
\def\mV{{\bm{V}}}
\def\mW{{\bm{W}}}
\def\mX{{\bm{X}}}
\def\mY{{\bm{Y}}}
\def\mZ{{\bm{Z}}}
\def\mBeta{{\bm{\beta}}}
\def\mPhi{{\bm{\Phi}}}
\def\mLambda{{\bm{\Lambda}}}
\def\mSigma{{\bm{\Sigma}}}

% Tensor
\DeclareMathAlphabet{\mathsfit}{\encodingdefault}{\sfdefault}{m}{sl}
\SetMathAlphabet{\mathsfit}{bold}{\encodingdefault}{\sfdefault}{bx}{n}
\newcommand{\tens}[1]{\bm{\mathsfit{#1}}}
\def\tA{{\tens{A}}}
\def\tB{{\tens{B}}}
\def\tC{{\tens{C}}}
\def\tD{{\tens{D}}}
\def\tE{{\tens{E}}}
\def\tF{{\tens{F}}}
\def\tG{{\tens{G}}}
\def\tH{{\tens{H}}}
\def\tI{{\tens{I}}}
\def\tJ{{\tens{J}}}
\def\tK{{\tens{K}}}
\def\tL{{\tens{L}}}
\def\tM{{\tens{M}}}
\def\tN{{\tens{N}}}
\def\tO{{\tens{O}}}
\def\tP{{\tens{P}}}
\def\tQ{{\tens{Q}}}
\def\tR{{\tens{R}}}
\def\tS{{\tens{S}}}
\def\tT{{\tens{T}}}
\def\tU{{\tens{U}}}
\def\tV{{\tens{V}}}
\def\tW{{\tens{W}}}
\def\tX{{\tens{X}}}
\def\tY{{\tens{Y}}}
\def\tZ{{\tens{Z}}}


% Graph
\def\gA{{\mathcal{A}}}
\def\gB{{\mathcal{B}}}
\def\gC{{\mathcal{C}}}
\def\gD{{\mathcal{D}}}
\def\gE{{\mathcal{E}}}
\def\gF{{\mathcal{F}}}
\def\gG{{\mathcal{G}}}
\def\gH{{\mathcal{H}}}
\def\gI{{\mathcal{I}}}
\def\gJ{{\mathcal{J}}}
\def\gK{{\mathcal{K}}}
\def\gL{{\mathcal{L}}}
\def\gM{{\mathcal{M}}}
\def\gN{{\mathcal{N}}}
\def\gO{{\mathcal{O}}}
\def\gP{{\mathcal{P}}}
\def\gQ{{\mathcal{Q}}}
\def\gR{{\mathcal{R}}}
\def\gS{{\mathcal{S}}}
\def\gT{{\mathcal{T}}}
\def\gU{{\mathcal{U}}}
\def\gV{{\mathcal{V}}}
\def\gW{{\mathcal{W}}}
\def\gX{{\mathcal{X}}}
\def\gY{{\mathcal{Y}}}
\def\gZ{{\mathcal{Z}}}

% Sets
\def\sA{{\mathbb{A}}}
\def\sB{{\mathbb{B}}}
\def\sC{{\mathbb{C}}}
\def\sD{{\mathbb{D}}}
% Don't use a set called E, because this would be the same as our symbol
% for expectation.
\def\sF{{\mathbb{F}}}
\def\sG{{\mathbb{G}}}
\def\sH{{\mathbb{H}}}
\def\sI{{\mathbb{I}}}
\def\sJ{{\mathbb{J}}}
\def\sK{{\mathbb{K}}}
\def\sL{{\mathbb{L}}}
\def\sM{{\mathbb{M}}}
\def\sN{{\mathbb{N}}}
\def\sO{{\mathbb{O}}}
\def\sP{{\mathbb{P}}}
\def\sQ{{\mathbb{Q}}}
\def\sR{{\mathbb{R}}}
\def\sS{{\mathbb{S}}}
\def\sT{{\mathbb{T}}}
\def\sU{{\mathbb{U}}}
\def\sV{{\mathbb{V}}}
\def\sW{{\mathbb{W}}}
\def\sX{{\mathbb{X}}}
\def\sY{{\mathbb{Y}}}
\def\sZ{{\mathbb{Z}}}

% Entries of a matrix
\def\emLambda{{\Lambda}}
\def\emA{{A}}
\def\emB{{B}}
\def\emC{{C}}
\def\emD{{D}}
\def\emE{{E}}
\def\emF{{F}}
\def\emG{{G}}
\def\emH{{H}}
\def\emI{{I}}
\def\emJ{{J}}
\def\emK{{K}}
\def\emL{{L}}
\def\emM{{M}}
\def\emN{{N}}
\def\emO{{O}}
\def\emP{{P}}
\def\emQ{{Q}}
\def\emR{{R}}
\def\emS{{S}}
\def\emT{{T}}
\def\emU{{U}}
\def\emV{{V}}
\def\emW{{W}}
\def\emX{{X}}
\def\emY{{Y}}
\def\emZ{{Z}}
\def\emSigma{{\Sigma}}

% entries of a tensor
% Same font as tensor, without \bm wrapper
\newcommand{\etens}[1]{\mathsfit{#1}}
\def\etLambda{{\etens{\Lambda}}}
\def\etA{{\etens{A}}}
\def\etB{{\etens{B}}}
\def\etC{{\etens{C}}}
\def\etD{{\etens{D}}}
\def\etE{{\etens{E}}}
\def\etF{{\etens{F}}}
\def\etG{{\etens{G}}}
\def\etH{{\etens{H}}}
\def\etI{{\etens{I}}}
\def\etJ{{\etens{J}}}
\def\etK{{\etens{K}}}
\def\etL{{\etens{L}}}
\def\etM{{\etens{M}}}
\def\etN{{\etens{N}}}
\def\etO{{\etens{O}}}
\def\etP{{\etens{P}}}
\def\etQ{{\etens{Q}}}
\def\etR{{\etens{R}}}
\def\etS{{\etens{S}}}
\def\etT{{\etens{T}}}
\def\etU{{\etens{U}}}
\def\etV{{\etens{V}}}
\def\etW{{\etens{W}}}
\def\etX{{\etens{X}}}
\def\etY{{\etens{Y}}}
\def\etZ{{\etens{Z}}}

% The true underlying data generating distribution
\newcommand{\pdata}{p_{\rm{data}}}
% The empirical distribution defined by the training set
\newcommand{\ptrain}{\hat{p}_{\rm{data}}}
\newcommand{\Ptrain}{\hat{P}_{\rm{data}}}
% The model distribution
\newcommand{\pmodel}{p_{\rm{model}}}
\newcommand{\Pmodel}{P_{\rm{model}}}
\newcommand{\ptildemodel}{\tilde{p}_{\rm{model}}}
% Stochastic autoencoder distributions
\newcommand{\pencode}{p_{\rm{encoder}}}
\newcommand{\pdecode}{p_{\rm{decoder}}}
\newcommand{\precons}{p_{\rm{reconstruct}}}

\newcommand{\laplace}{\mathrm{Laplace}} % Laplace distribution

\newcommand{\E}{\mathbb{E}}
\newcommand{\Ls}{\mathcal{L}}
\newcommand{\R}{\mathbb{R}}
\newcommand{\emp}{\tilde{p}}
\newcommand{\lr}{\alpha}
\newcommand{\reg}{\lambda}
\newcommand{\rect}{\mathrm{rectifier}}
\newcommand{\softmax}{\mathrm{softmax}}
\newcommand{\sigmoid}{\sigma}
\newcommand{\softplus}{\zeta}
\newcommand{\KL}{D_{\mathrm{KL}}}
\newcommand{\Var}{\mathrm{Var}}
\newcommand{\standarderror}{\mathrm{SE}}
\newcommand{\Cov}{\mathrm{Cov}}
% Wolfram Mathworld says $L^2$ is for function spaces and $\ell^2$ is for vectors
% But then they seem to use $L^2$ for vectors throughout the site, and so does
% wikipedia.
\newcommand{\normlzero}{L^0}
\newcommand{\normlone}{L^1}
\newcommand{\normltwo}{L^2}
\newcommand{\normlp}{L^p}
\newcommand{\normmax}{L^\infty}

\newcommand{\parents}{Pa} % See usage in notation.tex. Chosen to match Daphne's book.

\DeclareMathOperator*{\argmax}{arg\,max}
\DeclareMathOperator*{\argmin}{arg\,min}

\DeclareMathOperator{\sign}{sign}
\DeclareMathOperator{\Tr}{Tr}
\let\ab\allowbreak


\usepackage{hyperref}
\usepackage{url}
\hypersetup{
  colorlinks=true,
  allcolors=blue
}
\usepackage{graphicx}
\usepackage{booktabs}
\usepackage{wrapfig}
\usepackage{xcolor}
\usepackage{pifont}
\usepackage{multirow}
\usepackage{enumitem}
\usepackage{placeins}

% Continuous preprint layout: flexible floats with compact, readable spacing.
\setcounter{topnumber}{3}
\setcounter{bottomnumber}{2}
\setcounter{totalnumber}{5}
\renewcommand{\topfraction}{0.9}
\renewcommand{\bottomfraction}{0.8}
\renewcommand{\textfraction}{0.1}
\renewcommand{\floatpagefraction}{0.8}
\setlength{\textfloatsep}{12pt plus 2pt minus 2pt}
\setlength{\floatsep}{10pt plus 2pt minus 2pt}
\setlength{\intextsep}{10pt plus 2pt minus 2pt}

\newcommand{\fitwidth}[2]{%
  \sbox0{#2}%
  \ifdim\wd0>#1 \resizebox{#1}{!}{\usebox0}%
  \else \usebox0\fi}
\newcommand{\cmark}{\textcolor{green!60!black}{\ding{51}}} % 绿色勾
\newcommand{\xmark}{\textcolor{red!70!black}{\ding{55}}}   % 红色叉


% \title{World2Motion: Turning Video World Models into Third-Person Human Motion Generators}
\title{World2Motion: Turning Video World Models into 3D Human Motion Generators}

% arXiv author order follows OpenReview CHkfm9onjM, with Kaipeng Zhang
% inserted after Zhixiang Wang as requested.

% Affiliations: World2Motion OpenReview page; Kunhang Li and Kaipeng Zhang
% use their OpenReview career profiles (checked 2026-09-29).
% YIYI CAI and Bo Zheng affiliations corrected by the author on 2026-09-29.
\author{\bfseries Fangyuan Tu$^{1}$, Xiangyue Zhang$^{2}$, YIYI CAI$^{2}$, YICHEN PENG$^{4}$ \\
\bfseries Kunhang Li$^{2}$, Bo Zheng$^{3}$, Zhixiang Wang$^{3}$, Kaipeng Zhang$^{3}$ \\
\bfseries Erwin Wu$^{4}$, Haoran Xie$^{1}$, Haiyang Liu$^{2}$ \\[4pt]
{\normalfont\small $^{1}$Japan Advanced Institute of Science and Technology} \\
{\normalfont\small $^{2}$The University of Tokyo \quad $^{3}$Alaya Lab} \\
{\normalfont\small $^{4}$Institute of Science Tokyo}
}

% The \author macro works with any number of authors. There are two commands
% used to separate the names and addresses of multiple authors: \And and \AND.
%
% Using \And between authors leaves it to \LaTeX{} to determine where to break
% the lines. Using \AND forces a linebreak at that point. So, if \LaTeX{}
% puts 3 of 4 authors names on the first line, and the last on the second
% line, try using \AND instead of \And before the third author name.

\newcommand{\fix}{\marginpar{FIX}}
\newcommand{\new}{\marginpar{NEW}}

% Preprint formatting: show authors and suppress review line numbers.
% The conference publication header is replaced immediately after \maketitle.
\iclrfinalcopy
\begin{document}


\maketitle
\fancyhead{}
\lhead{Preprint}

\begin{abstract}
\begin{figure}[h]
    \centering
    \includegraphics[width=\linewidth]{Fig/pdf/Teaser_v3.pdf}
    % \caption{We propose World2Motion, turning video world models to synthesize realistic videos and 3d human motions synchronously based on a Mixture-of-Transformer(MoT) architecture. The top row shows the generated video frames, while the bottom row shows the generated human motions rendered in a 3D scene.}
    \caption{World2Motion adapts a video world model to jointly generate scene-aware 3D human motion and corresponding video from a single image and a text prompt. The example shows a person approaching a stool, turning, and sitting down. The top row shows generated video frames, while the bottom row shows the corresponding 3D motion rendered in a scene for visualization.}
    \label{fig:teaser}
\end{figure}

% Previous abstract preserved.
% Text-to-Motion(T2M) methods generate human motions efficiently but lack scene-aware capability. Human Object Interaction(HOI) methods can synthesize interactions, whereas they requires additional 3D object information. Video world models capture rich priors about how humans interact with the surrounding world. We aim to synthesize diverse human motions by exploiting the understanding of the physical world in video generation models. A straightforward solution is to generate various videos from prompts and extracted motions by Human Motion Recovery(HMR) methods. However, this two-stage paradigm is computationally expensive. To address these challenges, We present World2Motion, which turns video world model to synthesize the motions and videos synchronously. Directly adding motion branch to the video models is not efficiency. Since the videos and motions have fundamentally different representation space, it's costly for video generation backbone to generate motions. To address this challenge, we leverage a Mixture-of-Transformer(MoT) backbone to synthesize the motions and videos by a shared attention. However, naively fine-tuning the MoT backbone will introduce severe jitter into the generated motions. To address this challenge, we leverage a coupled noise schedule to alleviate the jitter. To train the MoT framework needs paired RGB videos and 3D human motions in interactive 3D scenes, existing HOI datasets can't satisfy the requirements. We curate a dataset containing 3D human motions and realistic videos. Comparison experiments demonstrate the efficiency, convenience and scene awareness of  our method. The ablation studies validate the effectiveness of our strategies.
% Before agreed claim-scope revision: We present World2Motion, a framework that generates scene-aware 3D human motion and corresponding video from a single image and a text prompt. While existing 3D motion generators learn from motion datasets, their generalization is constrained by limited coverage of environments. In contrast, video world models such as Cosmos 3 offer broader environmental priors but are not designed for full-body motion generation; recovering motion from their generated videos requires costly two-stage inference. To address these, we turn Cosmos 3 into a single-stage 3D motion generator. This adaptation has two challenges: the scarcity of paired video--motion data and temporal instability in the generated motion. First, we construct a training dataset combining synthetic video--motion pairs with real videos paired with estimated 3D motion. Second, we propose a shift-decoupled noise schedule that assigns different noise levels to video and motion through shared denoising progress. This design accommodates the different denoising requirements of the two modalities, reducing motion jitter. Experiments on a diverse test set assembled from multiple data sources show that World2Motion has better motion--text alignment and scene interaction compared with the evaluated 3D motion generators. It also matches the performance of the two-stage baseline while achieving approximately 3.3$\times$ faster inference.
We present World2Motion, a framework that generates scene-aware 3D human motion and corresponding video from a single image and a text prompt. While existing 3D motion generators learn from motion datasets, their generalization is constrained by limited coverage of environments. In contrast, video world models such as Cosmos 3 offer broader environmental priors but are not designed for full-body motion generation; recovering motion from their generated videos requires costly two-stage inference. To address these, we turn Cosmos 3 into a single-stage 3D motion generator. This adaptation has two challenges: the scarcity of paired video--motion data and temporal instability in the generated motion. First, we construct a training dataset combining synthetic video--motion pairs with real videos paired with estimated 3D motion. Second, we propose a shift-decoupled noise schedule that assigns different noise levels to video and motion through shared denoising progress. This design accommodates the different denoising requirements of the two modalities, reducing motion jitter. Experiments on a multi-source interaction benchmark show that World2Motion has better motion--text alignment and scene interaction compared with the evaluated 3D motion generators. It also matches the interaction success rate of the two-stage baseline while achieving approximately 3.3$\times$ faster inference. Our project page is available at \url{https://fyantu.github.io/World2Motion/}. 

% 1. leverage MoT architecture to solve the challenge 
% 2. dual noise schedule
% 3. dataset and experiment

\end{abstract}


% Main text: each section lives in its own file under sec/
% Previous introduction preserved (2026-09-26).
% \section{Introduction}
% \label{sec:intro}
% The generation of plausible 3D human motion is fundamental to modeling and understanding human behavior. Physically realistic motion sequences can benefit broad downstream applications, such as gaming and humanoid robotics. Conventional Text-to-Motion(T2M) methods~\citep{zhang2024motiondiffuse,tevet2023mdm} rapidly synthesize human motions from textual descriptions. Motivated by the scaling behavior observed in large language models, several approaches~\citep{hymotion2025, Kimodo2026} have trained on large-scale motion datasets to improve generation quality and generalization. T2M methods heavily rely on the text prompts to obtain the human motions. However, pure language struggle to describe detailed scene context, which limiting the ability to perceive and interact with the world. For example, it is difficult for T2M methods to generate a motion that precisely sits on a chair at a specified location. Achieving such interaction often requires manual adjustment of pose keyframes, which is cumbersome. Human-Object Interaction(HOI) methods~\citep{li2024chois, jung2026dechoi} are able to directly generate interaction motions from prompts. Although these methods perceive the environments, they usually rely on explicit 3D object information, such as object meshes ~\citep{kulkarni2024nifty} or object stats~\citep{li2024chois, zhang2026giraf}. Such strict data requirements make large-scale HOI data collection difficult, limiting these methods to leverage the data priors available in large motion data for further applications.
% 
% Video generation models have achieved remarkable progress in recent years. These models~\citep{wan2025wan, minimaxh3} produce high-quality human videos with strong visual fidelity and temporal consistency, demonstrating a strong capability to model the physical world and understand the human-environment interactions. With the development of deep learning, existing human mesh recovery(HMR)~\citep{patel2025camerahmr, zhao2026onlinehmr} methods have achieved strong performance in extracting continuous human motions from input images. Therefore, a straightforward method to generate plausible motions can be formulated as a two-stage pipeline: generate high-quality videos from text prompts and then extract the human motions using an HMR method. Although this approach enables scene-aware motion generations, its two-stage computational process is computationally expensive and time-consuming. 
% 
% \newsavebox{\introtabbox}
% \sbox{\introtabbox}{%
% \begin{tabular}{cccc}
% \toprule
%     & Speed & Convenience & scene awareness \\
% \midrule
% T2M     & \cmark    & \cmark    & \xmark    \\
% I2V+HMR & \xmark    & \cmark    & \cmark    \\
% HOI     & \cmark    & \xmark    & \cmark    \\
% Ours    & \cmark    & \cmark    & \cmark    \\
% \bottomrule
% \end{tabular}}
% \begin{wraptable}{r}{\wd\introtabbox}
% \centering
% \caption{Comparison of Motion Generation Paradigms.}
% \label{tab:paradigm comparison}
% \usebox{\introtabbox}
% \end{wraptable}
% 
% Video generation models incorporate rich visual priors about human motions and interactions in the world. Therefore, we try to leverage these priors to generate scene-aware human motions with plausible environmental interactions. To achieve these goal, we introduce World2Motion, which turns video world model to generates realistic videos and scene-aware 3D human motions synchronously. This joint generation paradigm allows the motion generation to benefit from the video generation models. Compared with the two-stage pipeline that first performs Image-to-video(I2V) generation and then recovers human motions using HMR approaches, our method achieves faster inference. Compared with conventional T2M method that rely solely on textual inputs, our framework additionally input an image to provide scene context. This enables the model to obtain more spatially accurate scene-aware 3D motions. Compared with existing HOI methods that require explicit 3D object geometry and object orientations, our method obtains the motions only from a single image. This paradigm makes it easier to input and more general to other scenarios that HOI methods can't covered, such as dancing. Overall, we characterize motion generation paradigms along three dimensions: Inference speed, convenience and scene awareness. As shown in Table~\ref{tab:paradigm comparison}, our paradigm simultaneously satisfies three criteria, providing a more balanced solution than existing motion generation paradigms.
% 
% Directly Adding
% Training our model requires paired data consisting of accurate 3D human motions interacting with objects in 3D scenes and corresponding RGB videos. However, existing motion datasets~\citep{li2023omomo, zhang2025motionx++} rarely provide both. To address this limitation, we develop a data pipeline to construct paired data from existing interaction datasets~\citep{li2023omomo, hassan2021samp, jiang2023chairs}. Directly fine-tune MoT~\citep{agarwal2026cosmos} backbone will lead to severe jitter in the generated motion sequences. To alleviate this phenomenon, we design two strategies to mitigate the jitter. First, we design a motion representation by anchoring root translation to the initial pose. To enable motion generation conditioned on the initial pose, we discard auxiliary features used in other methods~\citep{zhao2026ardy, tevet2023mdm}, such as root velocity and angular velocity. We use 6D joint rotations to encode orientations instead. Second, we introduce a coupled noise schedule that constrains the training trajectories of different modalities to lie on a shared curve. This strategy aligns the noise schedule during training and inference, which improves the motion quality. Extensive ablation studies demonstrate the effectiveness of our design. In summary, our contributions are as follows: 
% \begin{itemize}[leftmargin=*]
% \item we propose a unified MoT architecture to generate motions and videos synchronously. This paradigm synthesizes scene-aware motions while maintaining fast inference and convenience.
% \item We propose a motion representation and a coupled modality schedule to reduce jitter in the generated motions. Extensive ablation studies demonstrate the effectiveness of both designs.
% \item We develop a data construction pipeline for paired training data and compare other methods with multiple metrics. Experiments demonstrate the effectiveness and robustness of our paradigm.
% \end{itemize}
% End previous introduction.

\section{Introduction}
\label{sec:intro}
Generating 3D human motion for interactive applications requires motions that follow text descriptions and remain consistent with the surrounding environment. Conventional text-to-motion (T2M) methods~\citep{zhang2024motiondiffuse,tevet2023mdm} synthesize human motion from textual descriptions, and recent approaches~\citep{hymotion2025,Kimodo2026} have improved generation quality by scaling up motion datasets. However, the limited coverage of environments and interactions in available 3D motion datasets constrains their generalization to diverse scenes. Text alone also makes it difficult to specify detailed spatial relationships. For example, generating a person sitting on a chair requires the motion to match the chair's position and orientation. Human--object interaction (HOI) methods~\citep{li2024chois,jung2026dechoi} incorporate such spatial information through explicit 3D inputs, such as object meshes~\citep{kulkarni2024nifty} or object states~\citep{li2024chois,zhang2026giraf}. However, collecting the corresponding motion and scene data at scale remains difficult.

Video world models offer an alternative source of knowledge about human motion and environmental interactions~\citep{wan2025wan,minimaxh3}. Trained on diverse videos, models such as Cosmos 3~\citep{agarwal2026cosmos} provide environmental priors that can complement the limited coverage of 3D motion datasets. However, they are not designed to directly produce the full-body 3D motion required by our task. A straightforward way to use these priors is to first generate a video and then recover motion using a human mesh recovery (HMR) model~\citep{patel2025camerahmr,zhao2026onlinehmr}. This two-stage pipeline enables scene-aware motion generation, but requires a separate motion-recovery pass over the generated video, increasing inference cost.

% Previous Table 1 preserved (2026-09-26).
% \newsavebox{\introtabbox}
% \sbox{\introtabbox}{%
% \begin{tabular}{cccc}
% \toprule
%     & Speed & Convenience & scene awareness \\
% \midrule
% T2M     & \cmark    & \cmark    & \xmark    \\
% I2V+HMR & \xmark    & \cmark    & \cmark    \\
% HOI     & \cmark    & \xmark    & \cmark    \\
% Ours    & \cmark    & \cmark    & \cmark    \\
% \bottomrule
% \end{tabular}}
% \begin{wraptable}{r}{\wd\introtabbox}
% \centering
% \caption{Comparison of Motion Generation Paradigms.}
% \label{tab:paradigm comparison}
% \usebox{\introtabbox}
% \end{wraptable}
% End previous Table 1.

% Previous full-width Table 1 preserved.
% \begin{table}[htbp]
% \centering
% \normalfont\color{black}
% \caption{Comparison of motion generation paradigms.}
% \label{tab:paradigm comparison}
% \begin{tabular*}{\linewidth}{@{\extracolsep{\fill}}lccc@{}}
% \toprule
% Approach & \shortstack{Direct 3D\\motion output} & \shortstack{No 3D\\object input} & \shortstack{Scene\\conditioning} \\
% \midrule
% T2M & \ding{51} & \ding{51} & \\
% I2V + HMR & & \ding{51} & \ding{51} \\
% HOI & \ding{51} & & \ding{51} \\
% World2Motion & \ding{51} & \ding{51} & \ding{51} \\
% \bottomrule
% \end{tabular*}
% \end{table}
% End previous full-width Table 1.

\begin{wraptable}{r}{0.50\textwidth}
\centering
\normalfont\small\color{black}
\setlength{\tabcolsep}{2pt}
\caption{Motion generation methods comparison.}
\label{tab:paradigm comparison}
\begin{tabular*}{\linewidth}{@{\extracolsep{\fill}}lccc@{}}
\toprule
Approach & \shortstack{Direct 3D\\motion\\output} & \shortstack{No 3D\\object\\input} & \shortstack{Scene\\conditioning} \\
\midrule
T2M & \ding{51} & \ding{51} & \\
I2V + HMR & & \ding{51} & \ding{51} \\
HOI & \ding{51} & & \ding{51} \\
World2Motion & \ding{51} & \ding{51} & \ding{51} \\
\bottomrule
\end{tabular*}
\end{wraptable}

% Before shortening (2026-09-26).
% To use these video priors directly for motion generation, we introduce World2Motion, which adapts Cosmos 3 into a single-stage framework that generates scene-aware 3D human motion and corresponding video from a single image and a text prompt. We build on its existing Mixture-of-Transformers (MoT) backbone~\citep{agarwal2026cosmos} and extend the generation process to include motion. This allows the model to generate 3D motion alongside video, avoiding a separate motion-recovery stage on the generated sequence. This adaptation presents two challenges: the scarcity of paired video--motion training data and temporal instability in the generated motion.
% Before restoring the Table 1 discussion (2026-09-26).
% We introduce World2Motion, which adapts the existing Mixture-of-Transformers (MoT) backbone of Cosmos 3~\citep{agarwal2026cosmos} to jointly generate scene-aware 3D human motion and corresponding video from a single image and a text prompt. This single-stage framework avoids a separate motion-recovery stage on the generated sequence. The adaptation presents two challenges: the scarcity of paired video--motion training data and temporal instability in the generated motion.
We introduce World2Motion, which adapts the existing Mixture-of-Transformers (MoT) backbone of Cosmos 3~\citep{agarwal2026cosmos} to jointly generate scene-aware 3D human motion and corresponding video from a single image and a text prompt. This single-stage framework avoids a separate motion-recovery stage on the generated sequence. The adaptation presents two challenges: the scarcity of paired video--motion training data and temporal instability in the generated motion. Table~\ref{tab:paradigm comparison} summarizes the input requirements and generation pipelines of these approaches.

% Before shortening (2026-09-26).
% To address the scarcity of paired supervision, we construct a hybrid training dataset combining synthetic video--motion pairs with real videos paired with estimated 3D motion. For the synthetic component, we use motion sequences from existing interaction datasets~\citep{li2023omomo,hassan2021samp,jiang2023chairs} to construct corresponding videos, providing paired supervision with known 3D motion. For the real component, we estimate 3D motion from real videos~\citep{patel2025camerahmr} to expand the coverage of environments and appearances. These two sources combine motion supervision from existing 3D datasets with the visual diversity of real videos, supporting the adaptation of the video world model to joint video--motion generation.
% Before restoring the data explanation (2026-09-26).
% To address the scarcity of paired supervision, we construct a hybrid training dataset combining synthetic video--motion pairs with real videos paired with estimated 3D motion. We construct synthetic videos from existing interaction datasets~\citep{li2023omomo,hassan2021samp,jiang2023chairs}, providing supervision with known 3D motion, and estimate motion from real videos~\citep{patel2025camerahmr} to expand the coverage of environments and appearances.
To address the scarcity of paired supervision, we construct a hybrid training dataset combining synthetic video--motion pairs with real videos paired with estimated 3D motion. For the synthetic component, we use motion sequences from existing interaction datasets~\citep{li2023omomo,hassan2021samp,jiang2023chairs} to construct corresponding videos, providing paired supervision with known 3D motion. For the real component, we estimate 3D motion from real videos~\citep{patel2025camerahmr} to expand the coverage of environments and appearances. These two sources combine motion supervision from existing 3D datasets with the visual diversity of real videos, supporting the adaptation of the video world model to joint video--motion generation.

% Before shortening (2026-09-26).
% Directly fine-tuning the backbone for joint video--motion generation produces noticeable motion jitter. To improve temporal stability, we employ a motion representation with root translations relative to the initial root position and 6D joint rotations~\citep{zhao2026ardy}. We further account for the different denoising requirements of video and motion through a shift-decoupled noise schedule. The proposed schedule uses modality-specific shifts to assign different noise levels to video and motion through shared denoising progress. We maintain the same relationship between video and motion noise levels during training and inference. Our ablation studies show that the shift-decoupled noise schedule reduces motion jitter.
Directly fine-tuning the backbone for joint video--motion generation produces noticeable motion jitter. To improve temporal stability, we employ root translations relative to the initial root position and 6D joint rotations~\citep{zhao2026ardy}. We propose a shift-decoupled noise schedule that assigns different noise levels to video and motion through shared denoising progress, accounting for their different denoising requirements. We maintain the same video--motion noise relationship during training and inference. Our ablation studies show that the shift-decoupled noise schedule reduces motion jitter.

% Before shortening (2026-09-26).
% We evaluate World2Motion on a diverse test set assembled from multiple data sources. Compared with the evaluated 3D motion generators, World2Motion improves motion--text alignment and scene interaction. It also achieves comparable interaction success to the evaluated two-stage video-generation-and-motion-recovery baseline while providing approximately 3.3$\times$ faster inference. Ablation studies further demonstrate the contributions of the motion representation and shift-decoupled noise schedule to temporal stability.
% Before agreed claim-scope revision: On a diverse test set assembled from multiple data sources, World2Motion improves motion--text alignment and scene interaction over the evaluated 3D motion generators. It also achieves comparable interaction success to the MiniMax-H3 + CameraHMR baseline with approximately 3.3$\times$ faster inference.
On a multi-source interaction benchmark, World2Motion improves motion--text alignment and scene interaction over the evaluated 3D motion generators. It also achieves comparable interaction success to the MiniMax-H3 + CameraHMR baseline with approximately 3.3$\times$ faster inference. 

Our contributions are as follows:
\begin{itemize}[leftmargin=*]
\item We adapt Cosmos 3 for single-stage joint generation of scene-aware 3D human motion and corresponding video from a single image and a text prompt, avoiding a separate motion-recovery stage on the generated video.
\item We construct a hybrid training dataset combining synthetic video--motion pairs with real videos paired with estimated 3D motion.
\item We propose a shift-decoupled noise schedule that uses modality-specific noise shifts with shared denoising progress and maintains the same video--motion noise relationship during training and inference, improving temporal stability.
\end{itemize}

% Previous related work preserved (2026-09-26).
% \section{Related Work}
% 
% \textbf{Text-driven Human Motion Synthesize.} Text-to-Motion(T2M) generation aims to syhthesize 3D human motion from natural language. It has been dominated by two main paradigms. Diffusion-based methods such as MotionDiffuse~\citep{zhang2024motiondiffuse} and MDM~\citep{tevet2023mdm} generate motions by iterative denoising process with a text prompt. Following work such as ReMoDiffuse~\citep{zhang2023remodiffuse} and MLD~\citep{chen2023mld} improves the motion fidelity via retrieve argumentation. Meanwhile, autoregressive methods discretize motion into tokens and generate them sequentially. T2M-GPT~\citep{zhang2023t2mgpt} generates discrete motion tokens autoregressively, while MotionGPT~\citep{jiang2023motiongpt} further connects motion tokens to large language models. More recently, masked generative modeling has emerged as a strong alternative. MoMask~\citep{guo2024momask} employs residual vector quantization to represent motion as multi-layer discrete tokens. BAMM~\citep{pinyoanuntapong2024bamm} combines masked modeling with autoregressive generation, while MaskControl~\citep{pinyoanuntapong2025maskcontrol} introduces spatial control into masked motion synthesis. Recent works scale T2M models with larger architectures and datasets. HY-Motion~\citep{hymotion2025} and kimodo~\citep{Kimodo2026} scale DiT-based flow matching to the billion parameters and leverages large-scale motion pretraining to improve the motion quality. However, these methods rely primarily on text prompts or manipulating pose keyframes for control, limiting their ability to generate scene-aware human motions.
% 
% \textbf{Image-based Human Pose Regression.} Recovering 3D human pose and shape from images is commonly formulated as human mesh recovery (HMR). Most HMR methods regress the parameters of a statistical body model such as SMPL~\citep{SMPL:2015}. HMR~\citep{kanazawa2018hmr} uses iterative error feedback and adversarial pose priors to establish the end-to-end regression paradigm. SPIN~\citep{kolotouros2019spin} further improves HMR through optimization-in-the-loop training. CLIFF~\citep{li2022cliff} leverages the human location in the image to improve the mesh recovery. HybrIK~\citep{li2021hybrik} combines analytical inverse kinematics with neural networks for more accurate pose recovery. Humans in 4D~\citep{goel2023humansin4d} adopts a transformer-based architecture to reconstruct humans and track them over time. Recent HMR methods have expanded HMR methods toward whole-body~\citep{cai2023smpler} and multi-person reconstruction~\citep{baradel2024multihmr}. CameraHMR~\citep{patel2025camerahmr} improves camera modeling for accurate human placement in camera space. MetricHMSR~\citep{song2026metrichmsr} recovers metric human mesh and 3D scene geometry using perspective camera cues. OnlineHMR~\citep{zhao2026onlinehmr} further extends HMR to causal video processing, enabling online recovery of world-grounded human motion. The motions generated by these methods are highly dependent on the input images or videos, limiting their ability to produce diverse motions. With the rapid progress of video generation, another motion generation method use powerful video generation model such as Minimax-H3~\citep{minimaxh3} to synthesize diverse human videos and then extracted various motions by HMR-based methods~\citep{patel2025camerahmr}. However, two-stage pipeline requires considerable inference time.
% 
% \textbf{Scene-aware Human Pose Prediction.} Human motions are affected by the surrounding environments. POSA~\citep{hassan2021posa} leveraged human-scene contact and semantic relationships for plausible human placement in 3D scenes. SAMP~\citep{hassan2021samp} extends static motion modeling to dynamic motion generation. COINS~\citep{zhao2022coins} further introduces semantic control for compositional human-scene interactions. HUMANISE~\citep{wang2022humanise} further incorporates language conditions to generate semantically consistent human motions in 3D scenes. With the development of generative models, Diffusion-based approaches have been applied to human-object interaction synthesize. SceneDiffuser~\citep{jiang2024scenediffuser} conditions physical and orientation constraints to generate scene-aware human motions. OMOMO~\citep{li2023omomo} directly produces dynamic human-object interactions from object trajectories, without explicit specifications of the target interactions. Move as You Say~\citep{wang2024move} leverages scene affordance as an intermediate representation to generate language-guided human motions in 3D environments.TRUMANS~\citep{jiang2024trumans} models long-term human-scene interactions as a diffusion-based autoregressive model conditioned on scene context and actions. InfBaGel~\citep{zou2026infbagel} employs an iterative refinement strategy to generate interaction motions with dynamic scene awareness. Although these methods generate plausible interactions, they heavily rely on the 3D input, which limits their applicability. Move-in-2d~\citep{huang2025movein2d} adapts the motions to the visual environment to directly synthesize human motion just from a scene image and a text prompt. However, it lacks explicit control over the initial human pose, which makes it difficult to synthesize motions. 
% 
% \label{sec:related}

% Allow the introduction to flow into related work in the preprint.
\section{Related Work}
\label{sec:related}

% Before restoring masked-model details (2026-09-26).
% \noindent Text-conditioned human motion generation. Text-to-motion (T2M) methods generate 3D human motion from natural language. Diffusion-based methods such as MotionDiffuse~\citep{zhang2024motiondiffuse} and MDM~\citep{tevet2023mdm} denoise motion sequences, while MLD~\citep{chen2023mld} performs diffusion in a learned motion latent space and ReMoDiffuse~\citep{zhang2023remodiffuse} incorporates retrieved motion examples. Autoregressive methods, including T2M-GPT~\citep{zhang2023t2mgpt} and MotionGPT~\citep{jiang2023motiongpt}, generate discrete motion tokens sequentially. Masked models such as MoMask~\citep{guo2024momask}, BAMM~\citep{pinyoanuntapong2024bamm}, and MaskControl~\citep{pinyoanuntapong2025maskcontrol} further explore token prediction for motion generation and control. Recent work has also studied larger training datasets and more flexible motion constraints. HY-Motion~\citep{hymotion2025} scales DiT-based flow matching to a billion parameters, while Kimodo~\citep{Kimodo2026} trains a motion diffusion model on large-scale optical motion capture data and supports pose and trajectory constraints. These methods improve motion generation and control, but do not directly use scene images to model spatial relationships between a person and the surrounding environment.
\noindent \textbf{Text-conditioned human motion generation}. Text-to-motion (T2M) methods generate 3D human motion from natural language. Diffusion-based methods such as MotionDiffuse~\citep{zhang2024motiondiffuse} and MDM~\citep{tevet2023mdm} denoise motion sequences, while MLD~\citep{chen2023mld} performs diffusion in a learned motion latent space and ReMoDiffuse~\citep{zhang2023remodiffuse} incorporates retrieved motion examples. Autoregressive methods, including T2M-GPT~\citep{zhang2023t2mgpt} and MotionGPT~\citep{jiang2023motiongpt}, generate discrete motion tokens sequentially. Masked models such as MoMask~\citep{guo2024momask}, BAMM~\citep{pinyoanuntapong2024bamm}, and MaskControl~\citep{pinyoanuntapong2025maskcontrol} further explore token prediction for motion generation and control. MoMask uses residual vector quantization to represent motion with multiple token layers, while MaskControl incorporates spatial constraints. Recent work has also studied larger training datasets and more flexible motion constraints. HY-Motion~\citep{hymotion2025} scales DiT-based flow matching to a billion parameters, while Kimodo~\citep{Kimodo2026} trains a motion diffusion model on large-scale optical motion capture data and supports pose and trajectory constraints. These methods improve motion generation and control, but do not directly use scene images to model spatial relationships between a person and the surrounding environment.

% Before SceneDiffuser citation correction: \noindent Scene-conditioned human motion generation. Scene-conditioned methods incorporate environmental information into human pose and motion generation. Methods such as POSA~\citep{hassan2021posa}, SAMP~\citep{hassan2021samp}, COINS~\citep{zhao2022coins}, and HUMANISE~\citep{wang2022humanise} use 3D scene information to model human placement, contact, or motion. Other approaches explore different forms of interaction conditions: OMOMO~\citep{li2023omomo} uses object trajectories, Move as You Say~\citep{wang2024move} uses scene affordances, and TRUMANS~\citep{jiang2024trumans} conditions motion generation on scene context and actions. SceneDiffuser~\citep{jiang2024scenediffuser} and InfBaGel~\citep{zou2026infbagel} further study scene-conditioned generation and refinement for human interactions. These methods explore how scene information can guide motion beyond textual descriptions. Move-in-2D~\citep{huang2025movein2d} generates motion from a background image and text, and uses the generated motion to guide subsequent video generation. Our work leverages the priors learned during video world model pretraining to improve motion generalization across diverse environments.
% Before agreed claim-scope revision: \noindent Scene-conditioned human motion generation. Scene-conditioned methods incorporate environmental information into human pose and motion generation. Methods such as POSA~\citep{hassan2021posa}, SAMP~\citep{hassan2021samp}, COINS~\citep{zhao2022coins}, and HUMANISE~\citep{wang2022humanise} use 3D scene information to model human placement, contact, or motion. Other approaches explore different forms of interaction conditions: OMOMO~\citep{li2023omomo} uses object trajectories, Move as You Say~\citep{wang2024move} uses scene affordances, and TRUMANS~\citep{jiang2024trumans} conditions motion generation on scene context and actions. SceneDiffuser~\citep{huang2023scenediffuser} and InfBaGel~\citep{zou2026infbagel} further study scene-conditioned generation and refinement for human interactions. These methods explore how scene information can guide motion beyond textual descriptions. Move-in-2D~\citep{huang2025movein2d} generates motion from a background image and text, and uses the generated motion to guide subsequent video generation. Our work leverages the priors learned during video world model pretraining to improve motion generalization across diverse environments.
\noindent \textbf{Scene-conditioned human motion generation.} Scene-conditioned methods incorporate environmental information into human pose and motion generation. Methods such as POSA~\citep{hassan2021posa}, SAMP~\citep{hassan2021samp}, COINS~\citep{zhao2022coins}, and HUMANISE~\citep{wang2022humanise} use 3D scene information to model human placement, contact, or motion. Other approaches explore different forms of interaction conditions: OMOMO~\citep{li2023omomo} uses object trajectories, Move as You Say~\citep{wang2024move} uses scene affordances, and TRUMANS~\citep{jiang2024trumans} conditions motion generation on scene context and actions. SceneDiffuser~\citep{huang2023scenediffuser} and InfBaGel~\citep{zou2026infbagel} further study scene-conditioned generation and refinement for human interactions. These methods explore how scene information can guide motion beyond textual descriptions. Move-in-2D~\citep{huang2025movein2d} generates motion from a background image and text, and uses the generated motion to guide subsequent video generation. Our work studies the adaptation of a pretrained video world model for image-conditioned joint video and 3D motion generation.

\noindent \textbf{Human motion recovery from images and videos.} Human motion recovery estimates 3D body pose and shape from observed images or videos, often using a parametric body model such as SMPL~\citep{SMPL:2015}. HMR~\citep{kanazawa2018hmr}, SPIN~\citep{kolotouros2019spin}, CLIFF~\citep{li2022cliff}, and HybrIK~\citep{li2021hybrik} improve recovery through learned pose priors, optimization during training, image-location cues, and kinematic constraints. More recent methods study transformer-based recovery, whole-body and multi-person reconstruction, and camera and temporal modeling, including Humans in 4D~\citep{goel2023humansin4d}, SMPLer-X~\citep{cai2023smpler}, Multi-HMR~\citep{baradel2024multihmr}, CameraHMR~\citep{patel2025camerahmr}, MetricHMSR~\citep{song2026metrichmsr}, and OnlineHMR~\citep{zhao2026onlinehmr}. In our framework, CameraHMR provides motion estimates for real training videos and the initial pose from the input image. Motion recovery can also be applied to generated videos, forming a two-stage motion-generation pipeline, but the additional recovery stage increases inference cost. In contrast, World2Motion jointly generates motion and video, achieving approximately $3.3\times$ faster inference, better motion--text alignment, and lower motion jitter than the MiniMax-H3~\citep{minimaxh3} + CameraHMR~\citep{patel2025camerahmr} baseline.

\noindent \textbf{Video world models and joint video--motion generation.} Large-scale video models, such as Wan~\citep{wan2025wan} and MiniMax-H3~\citep{minimaxh3}, learn visual priors about human motion and environmental interactions. Cosmos 3~\citep{agarwal2026cosmos} extends multimodal generation to video and action within a Mixture-of-Transformers architecture. Its action modeling includes egocentric head and hand motion, and its action training inherits the vision noise schedule. This provides a starting point for adapting video world models to full-body 3D motion generation. Related work has also explored joint video--motion generation. CoMoVi~\citep{zhao2026comovi} uses separate diffusion branches for RGB video and 2D motion representations, together with 3D--2D cross-attention to predict 3D motion within the same denoising loop. World2Motion adapts the existing multimodal backbone of Cosmos 3 for joint video--motion generation. We combine synthetic and real video--motion training pairs and propose a shift-decoupled noise schedule that assigns modality-specific noise levels through shared denoising progress to improve temporal stability.

% Keep section transitions continuous in the preprint.
% Let floats share the following text page.
% Previous method preserved (2026-09-26).
% \section{Method}
% The understanding branch encodes the text into multi-level semantic representations, while the generation branch leverages the semantic features to generate the videos and human motions
% 
% Given an image $I_0 \in \mathbb{R}^{H \times W \times 3}$ containing a person in the scene, the corresponding human pose $P_0 \in \mathbb{R}^{J \times 3}$ and a text prompt of intended human actions, World2Motion jointly generates a video sequence $I_{1:N} \in \mathbb{R}^{N \times H \times W \times 3}$ and a human pose sequence $P_{1:N} \in \mathbb{R}^{N \times J \times 3}$ synchronously. Here, $H$ and $W$ denote the image height and width, $J$ denotes the number of SMPL~\citep{SMPL:2015} joints, and $N$ denotes the number of the generated sequence. Our framework adopts a Mixture-of-Transformers(MoT) architecture consisting of a understanding branch and a generation branch, as shown in Figure\ref{fig:Method}. 
% 
% % The text prompt is encoded into the understanding branch, while the video and the motion are processed by their respective encoders and fed into the generation branch. A shared attention layer is used to fuse information across multiple modalities. Finally, the video and motion outputs are decoded by respective decoders.
% 
% \begin{figure}[h]
%     \centering
%     \includegraphics[width=\linewidth]{Fig/pdf/Method_v3.pdf}
%     \caption{The architecture of our method.}
%     \label{fig:Method}
% \end{figure}
% 
% \subsection{Modality Tokenization}
% We first process the language, visual and motion inputs with modality-specific encoders to obtain the representations in a shared latent space. For text inputs, we adopt a text tokenizer to encode the prompts into the text embeddings. The text embeddings are provided as conditional context to make the generative branch better capture the semantic information of the prompt. To distinguish the text conditions from the generated modalities, we apply a 1D positional encoding to the text embeddings. The $i$-th text token $\mathbf{h}_{i}^{text}$ can be represented as:
% \begin{equation}
% \mathbf{h}_i^{\text{text}} = \mathbf{E}[w_i] \otimes \mathrm{RoPE}\big(\mathbf{p}_i^{\text{text}}\big),
% \qquad \mathbf{p}_i^{\text{text}} = (i,i,i)
% \label{eq:text-token}
% \end{equation}
% where $\mathbf{E}$ denotes the text embedding table, $w_i$ denotes the indices of the text embeddings, $\mathbf{p}_{i}^{text}$ denotes the coordinate of the $i$-th text token. We use a video VAE to compress the input videos into video latents. The video VAE applies $4\times$ temporal downsampling and $16\times$ spatial downsampling, followed by patchification with a $2\times2$ kernel. The video latent $x_{l,r,c}$ is projected by a vae2llm $W_v$ to align the dimension with the Transformer hidden space. A video-specified timestep emebedding $\boldsymbol{\phi}(\sigma_v)$ is then added to indicate the video noise level, followed by the video positional embedding $\mathrm{RoPE} \big(\mathbf{p}_{l,r,c}^{\text{video}}\big)$ to obtain the final video token $\mathbf{h}_{l,r,c}^{\text{video}}$:
% \begin{equation}
% \mathbf{h}_{l,r,c}^{\text{video}} = \Big( W_v x_{l,r,c} + \boldsymbol{\phi}(\sigma_v) \Big) \otimes \mathrm{RoPE} \big(\mathbf{p}_{l,r,c}^{\text{video}}\big), \qquad \mathbf{p}_{l,r,c}^{\text{video}}= (t_0+l, r, c)
% \label{eq:video-token}
% \end{equation}
% where $t_0, l, r, c$ indicate the origin of the temporal coordinate, video latent index, patch row index and patch column index respectively. $\sigma_v$ indicates the noise level of the video latent, and $\boldsymbol{\phi}(\cdot)$ represents the timestep embedding. 
% 
% \textbf{Motion Representation.} The input human motion $m_k$ of $k$-th frame contains root global orientation $r_0^k \in \mathbb{R}^3$, root global translation $\delta^k \in \mathbb{R}^3$ and joints orientations $r_{j_1, j_2,..j_n}^k \in \mathbb{R}^{n \times 3}$, $n$ represents the number of human joints. Then we convert the 3D axis-angle representation $r^k=[r_0^k, r_{j_1}^k,\ldots,r_{j_n}^k]$ into a 6D rotation representation $\tau^k=[\tau_0^k, \tau_{j_1}^k,\ldots,\tau_{j_n}^k]$ for more stable learning. Conventional Text-to-Motion (T2M) methods~\citep{guo2022humanml3d,zhang2023t2mgpt} generate an entire motion sequence directly from text. Their motion representations rely on root angular velocity and yaw increments to describe global motion and orientation. For our task, motion generation is conditioned on a initial pose without additional velocity. Therefore, we only use the 6D rotations and root translation to represent the pose. Directly feeding $\tau_k$ and $\delta_k$ into the network leads to temporal jitter in the generated motion. Inspired from ARDy~\citep{zhao2026ardy}, the root translation $\hat \delta_k$ is represented relative to the first frame translation $\delta_0$ as:
% \begin{equation}
% \hat \delta_k = \delta_k - \delta_0
% \label{eq:relative-translation}
% \end{equation}
% Then we apply the normalization and input the motion latent $x_k = [ \hat \delta_k, \tau_k ]$ into the encoder $W_a$ for dimensional alignment. We add a learnable embedding vector $e_{motion}$ to distinguish different modalities and use an action-specified timestep embedding $\boldsymbol{\phi}(\sigma_m)$ to indicate the motion noise level. Since the video latent is temporally $4 \times $ downsampled, we adjust the motion positional coordinate $\mathbf{p}_k^{\text{motion}}$ to align the video latent. Finally, the motion token $\mathbf{h}_k^{\text{motion}}$ can be described as:
% \begin{equation}
% \mathbf{h}_{k}^{\text{motion}} = \Big( W_a x_{k} + \mathbf{e}_{\text{motion}} + \boldsymbol{\phi}(\sigma_m) \Big) \otimes \mathrm{RoPE}\big(\mathbf{p}_k^{\text{motion}} \big), \qquad \mathbf{p}_k^{\text{motion}} = (t_0{+}\tfrac{k}{4}, 0, 0)
% \label{eq:motion-token}
% \end{equation}
% where $k$ denotes the index of the frame, $\sigma_m$ indicates the noise level of motion latent. $\mathrm{RoPE}(\cdot)$ denotes the rotary transform applied to the query and key projections of the token at every attention layer.
% 
% \subsection{Mixture of Transformer(MoT)}
% 
% Our backbone follows the MoT architecture design of Cosmos3~\citep{agarwal2026cosmos}, as illustrated in Figure\ref{fig:Method}(a). The MoT consists of an understanding branch and a generation branch. Each branch is a standard Transformer layer with its own parameters, including layer normalization modules, attention projection matrices, and feed-forward networks. The understanding branch implements a causal attention mask to achieve autoregressive design. During training, the understanding branch remains frozen and provides multi-level semantic guidance to the generation branch. The generation branch adopts a diffusion-based design, where the video and motion tokens are jointly fed into the Transformer for denoising. The understanding branch and generation branch interact through a shared multimodal attention layer. The Understanding branch attends causally within text tokens, while the generation branch attends bidirectionally to both AR and diffusion tokens through concatenated key and value representtaion.
% 
% \subsection{Coupled Noise Schedule}
% \textbf{Training.} A single noise level for both modalities is suboptimal: the video stream carries $\sim$25K tokens per sample while the motion stream carries around one hundred tokens. Different from EgoTwin~\citep{xiu2026egotwin}, we assign the two modalities different noise level through a shared variant $u$. This coupled scheme constrains all training noise pairs to lie on a single curve, ensuring better alignment with the inference trajectory. At each training step, we randomly sample a variant $u = g(\tau)$ with $\tau \sim \mathcal{U}(0,1)$ and warp it with per-modality shifts:
% 
% \begin{equation}
% \sigma_v = \frac{s_vu}{1+(s_v-1)u}, \qquad \sigma_m = \frac{s_mu}{1+(s_m-1)u}
% \label{eq:coupled-noise-schedule}
% \end{equation}
% 
% where $s_m$ and $s_v$ represent the shift of motion and video, respectively. We set $s_v$ higher than $s_m$ so that the video latent remains longer in the high-noise regime, which helps improve video generation quality. Both modality output are trained with rectified flow. Given a clean modality latent $z$, we construct noisy latent $z_\sigma = (1 - \sigma) z + \sigma\epsilon$. Each modality token receives its own timestep embedding, $\phi(\sigma_v)$ on video tokens and $\phi(\sigma_m)$ on motion tokens. Intuitively, the motion tokens is resolved earlier during denoising to provide the coarse structural guidance, while the video tokens refines high dimensional visual details. Finally, our training objective is formulated as:
% \begin{equation}
% \mathcal{L} = \mathcal{L}_{\text{video}} + \lambda\mathcal{L}_{\text{motion}}
% \label{eq:training-objective}
% \end{equation}
% where $\lambda$ is a coefficient for the motion loss, which compensates for the imbalance in token counts between the video and motion modalities.
% 
% \textbf{Inference.} Sampling runs two schedulers that share a single grid of progress values $u_{k}$ and differ only in their shifts $s_v, s_m$ at step $k$. At each inference step, the video and motion streams are assigned noise levels $\big(\sigma_v(u_k), \sigma_m(u_k)\big)$ along the same coupled trajectory used during training, ensuring that both modality latents denoise completely at the final step. After synchronous generation, we inverse-normalize the motion output and recover the global translation by adding the predicted relative translation to the positions of first frame. Finally, we apply a zero-phase Savitzky-Golay filter~\citep{savitzky1964smoothing} to the generated motion to improve the motion quality.
% 
% \begin{figure}[!h]
%     \centering
%     \includegraphics[width=\linewidth]{Fig/pdf/Denoising.pdf}
%     \caption{Qualitative comparison of motion generation methods.}
%     \label{fig:Comparison}
% \end{figure}
% 
% \label{sec:method}
% 
% End previous method.

\section{Method}
\label{sec:method}

Given a single image containing a person and a text prompt describing the intended action, World2Motion jointly generates a video sequence and corresponding 3D human motion. We estimate the initial human pose from the input image using CameraHMR~\citep{patel2025camerahmr}. To use the priors learned during video world model pretraining, we adapt the existing Cosmos 3 backbone for full-body motion generation. Figure~\ref{fig:Method} provides an overview of the framework.

\begin{figure}[!htbp]
    \centering
    \includegraphics[width=\linewidth]{Fig/pdf/Method_v6.pdf}
%     \caption{Overview of World2Motion. We adapt the Cosmos 3 backbone to jointly generate video and 3D human motion from an input image and a text prompt. Motion input and output projections connect the full-body motion representation to the pretrained backbone. A shift-decoupled noise schedule assigns modality-specific noise levels through shared denoising progress.}
%     \caption{Overview of World2Motion. Given an image and a text prompt, World2Motion jointly generates 3D human motion and video. CameraHMR estimates the initial pose from the image. A motion encoder and decoder adapt the Cosmos 3 backbone for full-body motion generation. The shift-decoupled noise schedule assigns modality-specific noise levels through shared denoising progress.}
    \caption{Overview of World2Motion. From an image and a text prompt, the adapted Cosmos 3 backbone jointly generates 3D human motion and video. CameraHMR estimates the initial pose from the image. The shift-decoupled schedule uses modality-specific noise shifts with shared denoising progress.}
    \label{fig:Method}
\end{figure}

\subsection{Review of Cosmos 3}
Cosmos 3~\citep{agarwal2026cosmos} uses a Mixture-of-Transformers (MoT) backbone with an understanding branch and a generation branch. We retain its text and video encoding components, shared attention mechanism, positional encoding, and train both generated modalities with rectified-flow objectives.

\subsection{Hybrid Video--Motion Training Data}
\label{sec:training-data}
% Before naming the dataset and specifying the real-video source: To address the scarcity of paired supervision, we construct a hybrid training dataset combining synthetic video--motion pairs with real videos paired with estimated 3D motion. For the synthetic component, we render motion sequences from existing interaction datasets, including OMOMO~\citep{li2023omomo}, CHAIRs~\citep{jiang2023chairs}, and SAMP~\citep{hassan2021samp}, and use the rendered animations to guide MiniMax-H3~\citep{minimaxh3} in generating corresponding RGB videos. The original motion sequences provide known 3D supervision. For the real component, we estimate 3D motion from internet videos using CameraHMR. We use Gemini 2.5~\citep{comanici2025gemini} to generate text descriptions for both components. These sources combine motion supervision from existing 3D datasets with the diversity of environments and appearances in real videos.
To address the scarcity of paired supervision, we construct HOI-mix, a hybrid dataset combining synthetic video--motion pairs with real videos paired with estimated 3D motion. For the synthetic component, we render motion sequences from existing interaction datasets, including OMOMO~\citep{li2023omomo}, CHAIRs~\citep{jiang2023chairs}, and SAMP~\citep{hassan2021samp}, and use the rendered animations to guide MiniMax-H3~\citep{minimaxh3} in generating corresponding RGB videos. The original motion sequences provide known 3D supervision. For the real component, we estimate 3D motion from curated Bilibili videos using CameraHMR. We use Gemini 2.5~\citep{comanici2025gemini} to generate text descriptions for both components. These sources combine motion supervision from existing 3D datasets with the diversity of environments and appearances in real videos.

\subsection{Full-Body Motion Adaptation}
% We represent each SMPL~\citep{SMPL:2015} motion frame using root translation and 6D joint rotations, including the root orientation. This representation allows generation to be conditioned on an initial pose without requiring velocity estimates. Following ARDY~\citep{zhao2026ardy}, we express root translation relative to its initial position:
% \begin{equation}
% \hat{\delta}_k = \delta_k - \delta_0,
% \label{eq:relative-translation}
% \end{equation}
% where $\delta_k$ is the global root translation at frame $k$ and $\delta_0$ is its initial position. We convert the axis-angle joint rotations to 6D rotations $\tau_k$, normalize the motion features $x_k=[\hat{\delta}_k,\tau_k]$, and project them into the Transformer hidden dimension using a motion encoder. A learnable modality embedding distinguishes motion tokens from video tokens. Motion time coordinates are scaled to match the video VAE's fourfold temporal compression, aligning the two modalities in time.
We represent each SMPL~\citep{SMPL:2015} motion frame using root translation and 6D joint rotations, including the root orientation. This representation allows generation to be conditioned on an initial pose without requiring velocity estimates. Following ARDY~\citep{zhao2026ardy}, we subtract the initial root position from the root translation at each frame. We convert the axis-angle joint rotations to 6D rotations, normalize the combined translation and rotation features, and project them into the Transformer hidden dimension using a motion encoder. A learnable modality embedding distinguishes motion tokens from video tokens. Motion time coordinates are scaled to match the video VAE's fourfold temporal compression, aligning the two modalities in time.

We train the motion encoder and decoder from scratch and initialize the remaining model weights from Cosmos 3. The understanding branch remains frozen during training. Our ablation study compares relative and absolute root translations and shows lower jerk and foot sliding with the relative representation (Table~\ref{tab:ablation_norm}).

\begin{figure}[!tb]
    \centering
    \includegraphics[width=\linewidth]{Fig/pdf/Denoising_v2.pdf}
%     \caption{Shift-decoupled video--motion denoising. (a) Video and motion noise levels over inference steps. (b) Shared denoising progress defines a coupled curve of video and motion noise levels. (c) Intermediate video frames and motion states at selected denoising steps. Rows show denoising steps, and column groups show different video times. The chair mesh is reference geometry for visualization.}
    \caption{Shift-decoupled video--motion denoising. (a) Video and motion noise levels during inference with $s_v=10$ and $s_m=1$. (b) The corresponding training noise pairs follow a shared curve, with lower motion noise than video noise at intermediate steps. (c) Intermediate video frames and motion states. Rows show denoising steps, and column groups show video times. The chair mesh is reference geometry for visualization.}
    \label{fig:denoising}
\end{figure}

\subsection{Shift-Decoupled Noise Schedule}
% Video and motion have different denoising requirements. We propose a shift-decoupled noise schedule that assigns modality-specific noise levels through a shared progress variable $u\in[0,1]$. At each training step, we sample one value of $u$ and compute the video and motion noise levels as
To improve temporal stability in joint video--motion generation, we propose a shift-decoupled noise schedule. It uses separate noise shifts for video and motion while retaining shared denoising progress $u\in[0,1]$. At each training step, we sample one value of $u$ and compute the video and motion noise levels as
\begin{equation}
\sigma_v(u) = \frac{s_vu}{1+(s_v-1)u}, \qquad
\sigma_m(u) = \frac{s_mu}{1+(s_m-1)u},
\label{eq:coupled-noise-schedule}
\end{equation}
% where $s_v$ and $s_m$ are positive shifts for video and motion, respectively. The shifts control how shared progress maps to each modality's noise level. When $s_v>s_m$, video has a higher noise level than motion for $0<u<1$. Each modality receives a timestep embedding corresponding to its own noise level. The shared progress constrains the training noise pairs to a single curve in the video--motion noise space, as illustrated in Figure~\ref{fig:denoising}(b).
where $s_v$ and $s_m$ are positive shifts for video and motion, respectively. The shifts control how shared progress maps to each modality's noise level. Setting $s_v>s_m$ keeps video at a higher noise level than motion during intermediate denoising steps, as illustrated in Figure~\ref{fig:denoising}(a). Each modality receives a timestep embedding corresponding to its own noise level. The shared progress constrains the training noise pairs to a single curve in the video--motion noise space, as illustrated in Figure~\ref{fig:denoising}(b).

% During inference, both modalities follow a common progress grid from $u=1$ to $u=0$ using the same shifts and mappings as in training. At each step, the video and motion streams use their respective noise levels along this shared curve. Both noise schedules reach zero at the final step. Our ablation studies show lower motion jitter with $s_v=10$ and $s_m=1$ than with equal shifts, and lower jitter when training and inference shifts match than when they differ (Table~\ref{tab:ablation_schedule}).
During inference, both modalities follow a common progress grid from $u=1$ to $u=0$, using the same shifts and mappings as in training. Both noise schedules reach zero at the final step. Among the tested configurations, $s_v=10$ and $s_m=1$ achieve the lowest jerk and foot sliding rate. Matching the inference shifts to the training shifts also improves both metrics compared with mismatched settings (Table~\ref{tab:ablation_schedule}).

After joint generation, we inverse-normalize the motion output and recover global root translation by adding the initial root position to the predicted relative translation. We then apply a Savitzky--Golay filter~\citep{savitzky1964smoothing} as post-processing. We disable this filter in all ablation experiments to evaluate the effects of the motion representation and noise schedule directly.



% Previous experiments preserved before nine-page revision (2026-09-26).
% \section{Experiments}
% 
% \subsection{Implement details}
% % text tokenizer
% We train the our model based on Cosmos3 NANO (16B) model~\citep{agarwal2026cosmos}, which consists of a Qwen3-VL(8B) text backbone~\citep{bai2025qwen3}, 36 MoT layers and a Wan2.2 causal VAE~\citep{wan2025wan}. The motion encoder and decoder are trained from scratch while all remaining weights are initialized from the base model. We train the network for 20,000 iterations using the AdamW optimizer with a learning rate of $5 \times 10^{-5}$, governed by a cosine decay schedule with 200 warmup steps. We randomly discard the text embeddings with a probability of 10$\%$ to achieve classifier-free guidance(CFG). Training is conducted in bfloat16(bf16) precision with a bath size of 16 across 16 NVIDIA H100 GPUs. The total training time is approximately 44h. 
% 
% \subsection{Dataset and Evaluation Metrics}
% Our dataset consists of two two parts: one is derived from 3D human-interaction dataset~\citep{li2023omomo, jiang2023chairs, hassan2021samp} and the other is collected from the internet. For the interaction datasets, we render the 3D human motions to obtain human animations. We then use the MiniMax-H3~\citep{minimaxh3} to generate the RGB videos from rendered animations. For the videos downloaded from the internet, we use cameraHMR~\citep{patel2025camerahmr} to obtain the human motions. All RGB videos are annotated by the gemini 2.5~\citep{comanici2025gemini}. Besides, we curate a test dataset from~\citet{li2023omomo, hassan2021samp}, COUCH~\citep{zhang2022couch} and internet to compare with other methods.
% 
% We evaluate the methods from multiply perspectives: 
% 
% \textbf{Motion Quality}: For Motion-prompt alignement, we follow the evaluation protocol defined by MoMask~\citep{guo2024momask} to evaluate the motion-prompt alignment using Frechet Inception Distance(FID), R-Precision(RP)  and MultiModal Distance (MM-Dist). For motion quality, we use Jerk and Foot Sliding Rate(FSR) followed from WHAM~\citep{shin2024wham} and OmniControl to evaluate the jitter~\citep{xie2024omnicontrol} respectively. Since we produce the videos and motions synchronously, we recover the motions from the videos using cameraHMR~\citep{patel2025camerahmr} and calculate the PoseError between extracted motions and generated motions to evaluate the consistency. 
% 
% \textbf{Scene Awareness}: Follow from OMOMO~\citep{li2023omomo}, we adopt contact-F1 as metrics. Besides, we render the generated motions in the corresponding 3D scenes and use QWen3-VL~\citep{bai2025qwen3} to determine whether the interaction is successfully completed. We define the ratio of scuccessful interaction cases over all test cases as the Interaction Success Rate(ISR).
% 
% \textbf{Speed}: We record the mean inference time of each method for Speed comparisons.
% 
% 
% \subsection{Comparison Results}
% We compare our method with three types of method: T2M methods, HOI methods and the I2V+HMR method. For T2M methods, we compare with Kimodo~\citep{Kimodo2026}, HY-Motion~\citep{hymotion2025} and Go-to-Zero~\citep{fan2025gotozero}. For HOI methods, we compare with LIGHT~\citep{wang2026light} and HOIFHLI~\citep{wu2025hoifhli}. For the two-stage inference method, we leverage MiniMax-H3~\citep{minimaxh3} to generate videos and CameraHMR~\citep{patel2025camerahmr} for subsequent motion recovery. CoMoVi~\citep{zhao2026comovi} is the most related work to ours. However, it additionally relies on a 2D motion representation for generation, whereas our method trains with less data requirements. we also incorporate CoMoVi~\citep{zhao2026comovi} as our baseline. For fair comparisons, the methods that conditioned on an initial human pose(Kimodo, HOIFHLI, CoMoVi, Ours) use the motion extracted by cameraHMR, which may introduce some space errors(Figure~\ref{fig:Comparison}(a)). In addition, we fine-tune CoMoVi on our training dataset for comparison. Since HOI methods generate human motions and object motions simultaneously, we will visualize the object trajectories generated by HOIFHLI and LIGHT. The complete quantitative results and qualitative results are shown in Table~\ref{tab:comparison} and Figure~\ref{fig:Comparison}.
% 
% \begin{table}[!htbp]
% \caption{Quantitative results of motion generation methods.}
% \label{tab:comparison}
% \begin{center}
% \resizebox{\textwidth}{!}{%
% \begin{tabular}{lccccccccc}
% \toprule
% \multirow{2}{*}{\textbf{Method}}
% & \multicolumn{3}{c}{\textbf{Motion-Prompt Alignment}}
% & \multicolumn{2}{c}{\textbf{Scene Awareness}}
% & \multicolumn{2}{c}{\textbf{Jitter}}
% & \multirow{2}{*}{\textbf{PoseError} $\downarrow$}
% & \multirow{2}{*}{\textbf{Time} $\downarrow$} \\
% \cmidrule(lr){2-4}
% \cmidrule(lr){5-6}
% \cmidrule(lr){7-8}
% & FID $\downarrow$
% & RP@1 $\uparrow$
% & MM-Dist $\downarrow$
% & Contact-F1 $\uparrow$
% & ISR $\uparrow$
% & Jerk $\downarrow$
% & FSR $\downarrow$
% & & \\
% 
% \midrule
% \textbf{T2M methods} & & & & & & & & & \\
% 
% \hspace{1em}HY-Motion
% & 0.46
% & 0.10
% & 1.21
% & 0.09
% & 0.51
% & 3.43
% & \textbf{1.25}
% & --
% & 2.83s \\
% 
% \hspace{1em}Kimodo
% & 0.66
% & 0.05
% & 1.30
% & 0.03
% & 0.23
% & 35.14
% & 6.07
% & --
% & \textbf{1.93s} \\
% 
% \hspace{1em}Go-to-Zero
% & 0.51
% & 0.04
% & 1.35
% & 0.07
% & 0.36
% & 17.50
% & 3.42
% & --
% & 2.57s \\
% 
% \addlinespace[0.5em]
% \multicolumn{10}{l}{\textbf{HOI methods}} \\
% 
% \hspace{1em}HOIFHLI
% & 0.82
% & 0.02
% & 1.37
% & 0.09
% & 0.12
% & 10.62
% & 6.28
% & --
% & 13.94s \\
% 
% \hspace{1em}LIGHT
% & 0.62
% & 0.03
% & 1.35
% & 0.17
% & 0.46
% & \textbf{1.37}
% & 11.60
% & --
% & 25.88s \\
% 
% \addlinespace[0.5em]
% \textbf{I2V+HMR} & & & & & & & & & \\
% 
% \begin{tabular}[c]{@{}l@{}}
% \hspace{1em}MiniMax-H3 +\\
% \hspace{1em}CameraHMR
% \end{tabular}
% & 0.69
% & 0.10
% & 1.32
% & \textbf{0.24}
% & \textbf{0.91}
% & 37.22
% & 6.41
% & --
% & 523.67s \\
% 
% \addlinespace[0.5em]
% \multicolumn{10}{l}{\textbf{Baseline}} \\
% 
% \hspace{1em}CoMoVi
% & 0.71
% & 0.12
% & 1.25
% & 0.09
% & 0.69
% & 16.10
% & 24.45
% & 81.28
% & 270.45s \\
% 
% \midrule
% 
% \hspace{1em}\textbf{Ours}
% & \textbf{0.27}
% & \textbf{0.24}
% & \textbf{1.13}
% & 0.22
% & \textbf{0.91}
% & 8.35
% & 2.80
% & \textbf{67.50}
% & 156.59s \\
% 
% \bottomrule
% \end{tabular}%
% }
% \end{center}
% \end{table}
% 
% \begin{figure}[!h]
%     \centering
%     \includegraphics[width=\linewidth]{Fig/pdf/Comparison.pdf}
%     \caption{Qualitative comparison of motion generation methods.}
%     \label{fig:Comparison}
% \end{figure}
% 
% \textbf{T2M methods.} Our method significantly outperforms T2M methods in both Motion-Prompt alignment and scene-awareness. Due to the lack of scene awareness, T2M methods fail to sit at the intended locations. T2M methods generally achieve lower jitter and FSR by encoding the velocity and acceleartion into motion representations. Since the motion stream contains fewer tokens than the video stream, T2M methods nearly achieve real-time inference.
% 
% \textbf{HOI methods.} Both HOIFHLI and LIGHT rely on 3D object information as input. LIGHT additionally requires human shape, while HOIFHLI requires the initial human pose. Our method generally outperforms the HOI methods in motion-prompt alignment and jitter. Since LIGHT and HOIFHLI are mainly optimized for dynamic interaction, Our method achieve better scene awareness as shown in Figure~\ref{fig:Comparison}(b)(c). Since the HOI methods don't need to denoise video latents, they are faster than our methods.
% 
% \textbf{I2V+HMR.} Our method outperforms the MiniMax-H3+CameraHMR pipeline across all evaluated metrics except Contact-F1. Because jointly generating introduces a trade-off in video quality, which further affect interaction accuracy in a little cases. In particular, the two-stage pipeline requires about 3$\times$ longer inference time.
% 
% \textbf{Baseline.} Our method also achieve better performance than the baseline across all metrics. CoMoVi predicts all joint rotations without root translations, making it effective for motions with limited displacement(Figure\ref{fig:Comparison}(a)). When the human needs to move, the motion accuracy drops significantly(Figure~\ref{fig:Comparison}(b)(c)).
% 
% 
% \subsection{Abalation Study}
% To verify that our strategies effectively reduce jitter, we conduct two ablation studies to evaluate the effectiveness. We use Jerk and Foot Sliding Rate as metrics to quantify the jitter. For a direct comparisons, we don't apply zero-phase Savitzky-Golay filter to any experimental settings.
% 
% \textbf{Motion Representation.} To validate the effectiveness of our motion representation, we use the absolute translation of the root node for comparison. We train the model with two motion representations while maintaining other settings fixed. The results are shown in Table\ref{tab:ablation_norm}. The experiment shows that relative translation representation outperforms the absolute translation representation on both metrics. Directly modeling absolute translations introduces noisy in the global coordinate space. The errors are magnified when the predictions are rescaled back to the original coordinate range, resulting in more severe temporal jitter.
% % \begin{table}[h]
% % \caption{Ablation study on motion representation.}
% % \label{tab:ablation_norm}
% % \begin{center}
% % \begin{tabular}{lccc}
% % \toprule
% %  & Jerk$\downarrow$  & Foot sliding rate $\downarrow$ \\
% % \midrule
% % % w/o relative translation & 4.355 $\pm$ 4.425 & 0.766 & \\
% % % w/ relative translation & \textbf{1.251 $\pm$ 0.973} & \textbf{0.457} & \\
% % w/o relative translation & 4.355 & 0.766 & \\
% % w/ relative translation & \textbf{1.251} & \textbf{0.457} & \\
% % \bottomrule
% % \end{tabular}
% % \end{center}
% % \end{table}
% \begin{table}[h]
% \caption{Root translation ablation without post-processing. Translations are either absolute or relative to the initial root position; all other settings are fixed. Lower values are better.}
% \label{tab:ablation_norm}
% \begin{center}
% \begin{tabular}{lcc}
% \toprule
%  & Jerk $\downarrow$  & Foot sliding rate $\downarrow$ \\
% \midrule
% % w/o relative translation & 4.355 $\pm$ 4.425 & 0.766 & \\
% % w/ relative translation & \textbf{1.251 $\pm$ 0.973} & \textbf{0.457} & \\
% Absolute root translation & 4.355 & 0.766 \\
% Relative root translation & 1.251 & 0.457 \\
% \bottomrule
% \end{tabular}
% \end{center}
% \end{table}
% 
% 
% % \textbf{Coupled Dual-timestep Schedule.}
% \noindent Shift-Decoupled Noise Schedule.
% % TODO: 关于去噪理论的解释，假定每个token包含的信息密度一样，那么数量越多的token越需要进行去噪在高频，从而让去噪的平均值和token数少的保持一致，这点应该要重新再写下
% 
% % Before correcting the denoising figure reference (2026-09-26).
% % To validate the effectiveness of the coupled noise schedule and train-inference consistency, we conduct nine ablation experiments as shown in Table~\ref{tab:ablation_schedule}. We train three shift configurations while keeping other parameters fixed. For each training shift configurations, the lowest jitter is observed when the training and inference shifts are kept consistent. Since the model is trained on noise pairs along a single coupled trajectory (Figure\ref{fig:Method} (c)), Using noise pairs deviated from the trajectory will lead to prediction errors and more severe jitter. Comparing the different shift configurations, we find that \{ $s_v$ = 10, $s_m$ = 1 \} achieve the best performance. Since the motion latents is more compact and less redundancy, the uniform denoising schedule is more sufficient than keeping the latent longer in the high noise region.
% % To validate the effectiveness of the coupled noise schedule and train-inference consistency, we conduct nine ablation experiments as shown in Table~\ref{tab:ablation_schedule}. We train three shift configurations while keeping other parameters fixed. For each training shift configurations, the lowest jitter is observed when the training and inference shifts are kept consistent. Since the model is trained on noise pairs along a single coupled trajectory (Figure~\ref{fig:denoising}(b)), Using noise pairs deviated from the trajectory will lead to prediction errors and more severe jitter. Comparing the different shift configurations, we find that \{ $s_v$ = 10, $s_m$ = 1 \} achieve the best performance. Since the motion latents is more compact and less redundancy, the uniform denoising schedule is more sufficient than keeping the latent longer in the high noise region.
% We evaluate three training shift configurations, each with three inference configurations, while keeping other parameters fixed (Table~\ref{tab:ablation_schedule}). For every training configuration, matching the inference shifts to the training shifts gives the lowest jerk and foot sliding rate. Among all tested configurations, matched training and inference shifts of $s_v=10$ and $s_m=1$ achieve the lowest values on both metrics. These results support using modality-specific shifts and maintaining the same noise relationship during training and inference.
% 
% 
% % \begin{table}[!htbp]
% % \caption{Ablation study on Coupled Dual-timestep Schedule.}
% % \label{tab:ablation_schedule}
% % \begin{center}
% % \begin{tabular}{ccccc}
% % \toprule
% % Training & Inference & Jerk$\downarrow$ & Foot Sliding ratio$\downarrow$ \\
% % \midrule
% % % \multirow{3}{*}{$s_v=10,s_m=1$}
% % %  & $s_v=10, s_m=1$ & \textbf{1.251 $\pm$ 0.97} & \textbf{0.457} \\
% % %  & $s_v=10, s_m=10$ & 8.003 $\pm$ 2.48 & 0.986 \\
% % %  & $s_v=10, s_m=15$ & 9.703 $\pm$ 3.47 & 1.000\\
% % % \midrule
% % % \multirow{3}{*}{$s_v=10,s_m=10$}
% % %  & $s_v=10, s_m=1$ & 7.021 $\pm$ 1.04 & 0.985\\
% % %  & $s_v=10, s_m=10$ & \textbf{1.635} $\pm$ \textbf{1.05} & \textbf{0.635} \\
% % %  & $s_v=10, s_m=15$ & 2.557 $\pm$ 0.79 & 0.830 \\
% % % \midrule
% % % \multirow{3}{*}{$s_v=10,s_m=15$}
% % %  & $s_v=10, s_m=1$ & 8.413 $\pm$ 1.01 & 0.981\\
% % %  & $s_v=10, s_m=10$ & 2.127 $\pm$ 1.55 & 0.700\\
% % %  & $s_v=10, s_m=15$ & \textbf{1.809} $\pm$ \textbf{1.28} & \textbf{0.649} \\
% % \multirow{3}{*}{$s_v=10,s_m=1$}
% %  & $s_v=10, s_m=1$ & \textbf{1.251} & \textbf{0.457} \\
% %  & $s_v=10, s_m=10$ & 8.003 & 0.986 \\
% %  & $s_v=10, s_m=15$ & 9.703 & 1.000\\
% % \midrule
% % \multirow{3}{*}{$s_v=10,s_m=10$}
% %  & $s_v=10, s_m=1$ & 7.021 & 0.985\\
% %  & $s_v=10, s_m=10$ & \textbf{1.635} & \textbf{0.635} \\
% %  & $s_v=10, s_m=15$ & 2.557 & 0.830 \\
% % \midrule
% % \multirow{3}{*}{$s_v=10,s_m=15$}
% %  & $s_v=10, s_m=1$ & 8.413 & 0.981\\
% %  & $s_v=10, s_m=10$ & 2.127 & 0.700\\
% %  & $s_v=10, s_m=15$ & \textbf{1.809} & \textbf{0.649} \\
% % \bottomrule
% % \end{tabular}
% % \end{center}
% % \end{table}
% \begin{table}[!htbp]
% \caption{Ablation of the shift-decoupled noise schedule without post-processing. Each group fixes the training shifts and varies the inference shifts. Matching shifts yields the lowest jerk and foot sliding rate within each group.}
% \label{tab:ablation_schedule}
% \begin{center}
% \begin{tabular}{cccc}
% \toprule
% Training shifts & Inference shifts & Jerk $\downarrow$ & Foot sliding rate $\downarrow$ \\
% \midrule
% % \multirow{3}{*}{$s_v=10,s_m=1$}
% %  & $s_v=10, s_m=1$ & 1.251 $\pm$ 0.97 & 0.457 \\
% %  & $s_v=10, s_m=10$ & 8.003 $\pm$ 2.48 & 0.986 \\
% %  & $s_v=10, s_m=15$ & 9.703 $\pm$ 3.47 & 1.000\\
% % \midrule
% % \multirow{3}{*}{$s_v=10,s_m=10$}
% %  & $s_v=10, s_m=1$ & 7.021 $\pm$ 1.04 & 0.985\\
% %  & $s_v=10, s_m=10$ & 1.635 $\pm$ 1.05 & 0.635 \\
% %  & $s_v=10, s_m=15$ & 2.557 $\pm$ 0.79 & 0.830 \\
% % \midrule
% % \multirow{3}{*}{$s_v=10,s_m=15$}
% %  & $s_v=10, s_m=1$ & 8.413 $\pm$ 1.01 & 0.981\\
% %  & $s_v=10, s_m=10$ & 2.127 $\pm$ 1.55 & 0.700\\
% %  & $s_v=10, s_m=15$ & 1.809 $\pm$ 1.28 & 0.649 \\
% \multirow{3}{*}{$s_v=10,s_m=1$}
%  & $s_v=10, s_m=1$ & 1.251 & 0.457 \\
%  & $s_v=10, s_m=10$ & 8.003 & 0.986 \\
%  & $s_v=10, s_m=15$ & 9.703 & 1.000\\
% \midrule
% \multirow{3}{*}{$s_v=10,s_m=10$}
%  & $s_v=10, s_m=1$ & 7.021 & 0.985\\
%  & $s_v=10, s_m=10$ & 1.635 & 0.635 \\
%  & $s_v=10, s_m=15$ & 2.557 & 0.830 \\
% \midrule
% \multirow{3}{*}{$s_v=10,s_m=15$}
%  & $s_v=10, s_m=1$ & 8.413 & 0.981\\
%  & $s_v=10, s_m=10$ & 2.127 & 0.700\\
%  & $s_v=10, s_m=15$ & 1.809 & 0.649 \\
% \bottomrule
% \end{tabular}
% \end{center}
% \end{table}
% 
% 
% 
% \label{sec:experiments}
% 
% 

% Previous nine-page experiments preserved before structure and table revision.
% \clearpage
% \section{Experiments}
% \label{sec:experiments}
% \subsection{Experimental Setup}
% \label{sec:experimental-setup}
% 
% \noindent Training. We initialize World2Motion from Cosmos 3 NANO (16B)~\citep{agarwal2026cosmos}, which includes a Qwen3-VL (8B) text backbone~\citep{bai2025qwen3}, 36 Mixture-of-Transformers layers, and a Wan2.2 causal VAE~\citep{wan2025wan}. The motion encoder and decoder are trained from scratch, and the understanding branch remains frozen. We train for 20,000 iterations with AdamW, a learning rate of $5\times10^{-5}$, a cosine decay schedule, and 200 warmup steps. Text conditioning is dropped with probability 0.1 for classifier-free guidance. Training uses bfloat16 precision, a batch size of 16, and 16 NVIDIA H100 GPUs, and takes approximately 44 hours.
% 
% \noindent Training and test data. The hybrid training data combines synthetic video--motion pairs derived from OMOMO~\citep{li2023omomo}, CHAIRs~\citep{jiang2023chairs}, and SAMP~\citep{hassan2021samp} with internet videos paired with CameraHMR~\citep{patel2025camerahmr} motion estimates. Section~\ref{sec:training-data} describes their construction. The synthetic component supplies known 3D motion, while the real component expands the range of environments and appearances. Text descriptions are generated using Gemini 2.5~\citep{comanici2025gemini}. We evaluate on a test set assembled from OMOMO, SAMP, COUCH~\citep{zhang2022couch}, and internet videos. This multi-source evaluation tests whether generated motions follow the requested action and its spatial relationship to the depicted environment.
% 
% \noindent Baselines and conditioning. We compare with text-to-motion methods HY-Motion~\citep{hymotion2025}, Kimodo~\citep{Kimodo2026}, and Go-to-Zero~\citep{fan2025gotozero}; human--object interaction methods LIGHT~\citep{wang2026light} and HOIFHLI~\citep{wu2025hoifhli}; and a two-stage MiniMax-H3~\citep{minimaxh3} + CameraHMR pipeline. We also fine-tune CoMoVi~\citep{zhao2026comovi} on our training data as a joint video--motion baseline. Methods that accept an initial human pose---Kimodo, HOIFHLI, CoMoVi, and World2Motion---use the pose estimated by CameraHMR. This shared pose estimate can contain spatial errors that affect the subsequent motion. LIGHT and HOIFHLI additionally require explicit 3D object information; World2Motion obtains scene context from the input image. For the interaction baselines, we visualize their predicted object trajectories together with the human motion.
% 
% \noindent Motion distribution and text alignment. We follow the MoMask~\citep{guo2024momask} evaluation protocol. Fr\'echet Inception Distance (FID) measures the difference between generated and reference motion feature distributions. R-Precision at rank 1 (R@1) and multimodal distance (MM-Dist) evaluate motion--text alignment. Lower FID and MM-Dist and higher R@1 are better. These scores assess different properties: a motion can match an action description while failing to reach the intended object.
% 
% \noindent Scene interaction and motion stability. Following OMOMO, we report Contact-F1. We also render generated motions in the corresponding 3D scenes and use Qwen3-VL to assess whether the requested interaction is completed. Interaction Success Rate (ISR) is the fraction of test cases judged successful. Jerk measures temporal motion irregularity, and foot sliding rate (FSR) measures sliding artifacts, following the evaluation procedures of WHAM~\citep{shin2024wham} and OmniControl~\citep{xie2024omnicontrol}. Lower values are better for both stability metrics.
% 
% \noindent Video--motion consistency and efficiency. For joint generation methods, we recover motion from the generated video using CameraHMR and report PoseError against the directly generated motion. This evaluates agreement between the two outputs. We also report mean inference time in seconds. World2Motion results in the main comparison use the post-processing described in Section~\ref{sec:method}; all ablation results are obtained without the Savitzky--Golay filter so that representation and scheduling choices can be assessed directly.
% 
% \noindent Interpreting the metrics. PoseError depends on video-based motion recovery and measures cross-modal agreement rather than error against a ground-truth 3D trajectory. We therefore read it alongside text alignment and interaction scores. Similarly, low jerk does not establish that a requested interaction is completed: smooth motion can miss an object, while a successful interaction may still contain local artifacts. Reporting these measures together separates scene interaction from temporal stability and provides a more complete assessment of scene-aware motion generation.
% 
% \clearpage
% \subsection{Quantitative Comparisons}
% 
% \begin{table}[!htbp]
% \centering
% \small
% \setlength{\tabcolsep}{4pt}
% \caption{Comparison on the multi-source test set. (a) Motion distribution, text alignment, and scene interaction. (b) Motion stability, video--motion consistency, and mean inference time. Arrows indicate better values; underlining marks the best reported value in each column. Dashes indicate unreported scores.}
% \label{tab:comparison}
% \begin{tabular*}{\linewidth}{@{\extracolsep{\fill}}lccccc@{}}
% \toprule
% \multicolumn{6}{l}{(a) Motion generation and scene interaction}\\
% Method & FID $\downarrow$ & R@1 $\uparrow$ & MM-Dist $\downarrow$ & Contact-F1 $\uparrow$ & ISR $\uparrow$\\
% \midrule
% \multicolumn{6}{l}{Text-to-motion}\\
% HY-Motion & 0.46 & 0.10 & 1.21 & 0.09 & 0.51\\
% Kimodo & 0.66 & 0.05 & 1.30 & 0.03 & 0.23\\
% Go-to-Zero & 0.51 & 0.04 & 1.35 & 0.07 & 0.36\\
% \midrule
% \multicolumn{6}{l}{Human--object interaction}\\
% HOIFHLI & 0.82 & 0.02 & 1.37 & 0.09 & 0.12\\
% LIGHT & 0.62 & 0.03 & 1.35 & 0.17 & 0.46\\
% \midrule
% \multicolumn{6}{l}{Two-stage video generation and motion recovery}\\
% MiniMax-H3 + CameraHMR & 0.69 & 0.10 & 1.32 & \underline{0.24} & \underline{0.91}\\
% \midrule
% \multicolumn{6}{l}{Joint video--motion generation}\\
% CoMoVi & 0.71 & 0.12 & 1.25 & 0.09 & 0.69\\
% World2Motion & \underline{0.27} & \underline{0.24} & \underline{1.13} & 0.22 & \underline{0.91}\\
% \bottomrule
% \end{tabular*}
% 
% \medskip
% \begin{tabular*}{\linewidth}{@{\extracolsep{\fill}}lcccc@{}}
% \toprule
% \multicolumn{5}{l}{(b) Motion stability, output consistency, and efficiency}\\
% Method & Jerk $\downarrow$ & FSR $\downarrow$ & PoseError $\downarrow$ & Time (s) $\downarrow$\\
% \midrule
% \multicolumn{5}{l}{Text-to-motion}\\
% HY-Motion & 3.43 & \underline{1.25} & -- & 2.83\\
% Kimodo & 35.14 & 6.07 & -- & \underline{1.93}\\
% Go-to-Zero & 17.50 & 3.42 & -- & 2.57\\
% \midrule
% \multicolumn{5}{l}{Human--object interaction}\\
% HOIFHLI & 10.62 & 6.28 & -- & 13.94\\
% LIGHT & \underline{1.37} & 11.60 & -- & 25.88\\
% \midrule
% \multicolumn{5}{l}{Two-stage video generation and motion recovery}\\
% MiniMax-H3 + CameraHMR & 37.22 & 6.41 & -- & 523.67\\
% \midrule
% \multicolumn{5}{l}{Joint video--motion generation}\\
% CoMoVi & 16.10 & 24.45 & 81.28 & 270.45\\
% World2Motion & 8.35 & 2.80 & \underline{67.50} & 156.59\\
% \bottomrule
% \end{tabular*}
% \end{table}
% 
% \noindent Motion generation and scene interaction. World2Motion achieves the lowest FID and MM-Dist and the highest R@1 among the compared methods (Table~\ref{tab:comparison}(a)). Relative to HY-Motion, R@1 increases from 0.10 to 0.24 and ISR from 0.51 to 0.91. Its ISR exceeds those of the evaluated text-to-motion and human--object interaction methods. The two-stage baseline reaches the same ISR and a higher Contact-F1 (0.24 versus 0.22). Thus, the results support improved action alignment and competitive scene interaction, while leaving room to improve contact accuracy.
% 
% \noindent Stability and output consistency. World2Motion has lower jerk and FSR than Kimodo, Go-to-Zero, HOIFHLI, and the two video-based baselines. HY-Motion achieves lower values on both stability metrics, and LIGHT achieves lower jerk. World2Motion therefore does not lead every measure of motion quality. Compared with CoMoVi, it reduces PoseError from 81.28 to 67.50, a 17.0\% reduction, indicating closer agreement between the generated video and 3D motion.
% 
% \noindent Inference efficiency. World2Motion takes 156.59 seconds on average, compared with 523.67 seconds for MiniMax-H3 + CameraHMR and 270.45 seconds for CoMoVi: speedups of approximately $3.34\times$ and $1.73\times$, respectively. The text-to-motion and human--object interaction methods remain faster. Our efficiency gain concerns video-based motion generation; jointly producing video and motion still costs more than generating motion alone.
% 
% \clearpage
% \subsection{Qualitative Comparisons}
% 
% \begin{figure}[!htbp]
% \centering
% \includegraphics[width=0.96\linewidth]{Fig/pdf/Comparison.pdf}
% \caption{Qualitative comparison of scene interactions: (a) lifting a plastic box, (b) walking to a chair and sitting, and (c) walking to a sofa and sitting. Each method shows successive motion states; video-generating methods also show corresponding video frames. Human--object interaction baselines include their predicted object trajectories.}
% \label{fig:Comparison}
% \end{figure}
% 
% Figure~\ref{fig:Comparison} compares actions that depend on both body motion and placement relative to an object. In the sitting examples, producing a plausible sitting pose is insufficient: the character must first approach the intended seat and then align the body with it. The box example also involves coordinating the body with object motion. Viewing the approach and final interaction together helps distinguish action matching from scene-consistent execution. For video-based methods, the paired rows additionally allow the reader to inspect whether the generated video and recovered or directly generated 3D motion describe the same interaction.
% 
% \clearpage
% \subsection{Ablation Studies}
% 
% We evaluate root translation and noise scheduling using jerk and FSR. All ablations disable post-processing. Comparisons within each study keep other settings fixed. The first study compares motion representations; the second changes training and inference shifts independently. This separates the choice of noise schedule from the effect of using it consistently.
% 
% \noindent Root translation. Representing each root position relative to the initial position reduces jerk from 4.355 to 1.251 (71.3\%) and FSR from 0.766 to 0.457 (40.3\%), compared with absolute translation (Table~\ref{tab:ablation_norm}). Both representations retain root position and 6D joint rotations; the change is the reference position used for translation. The comparison therefore isolates this coordinate choice while keeping the rotation representation and other settings fixed. These results support using relative positions for more stable generation without requiring root velocity as an input.
% 
% \begin{table}[!htbp]
% \centering
% \small
% \setlength{\tabcolsep}{8pt}
% \caption{Root translation ablation without post-processing. Relative positions are measured from the initial root position. Underlining marks the lower value in each column.}
% \label{tab:ablation_norm}
% \begin{tabular}{lcc}
% \toprule
% Root translation & Jerk $\downarrow$ & FSR $\downarrow$\\
% \midrule
% Absolute & 4.355 & 0.766\\
% Relative to initial position & \underline{1.251} & \underline{0.457}\\
% \bottomrule
% \end{tabular}
% \end{table}
% 
% \noindent Shift-decoupled noise schedule. Table~\ref{tab:ablation_schedule} separates two questions: which shifts to use, and whether the training and inference shifts should match. With matched shifts, $s_m=1$ achieves lower jerk and FSR than $s_m=10$ or $15$. Compared with matched $s_m=10$, matched $s_m=1$ reduces jerk by 23.5\% and FSR by 28.0\%. For each training configuration, matching the inference shift also gives the lowest values on both metrics. For example, after training with $s_m=1$, changing only the inference shift to $10$ increases jerk from 1.251 to 8.003. These observations support modality-specific shifts and consistent training and inference noise relationships. They do not require an additional smoothing step, since the filter is disabled throughout this comparison.
% 
% \begin{table}[!htbp]
% \centering
% \small
% \setlength{\tabcolsep}{10pt}
% \caption{Noise-schedule ablation without post-processing. Video shift $s_v=10$ is fixed in training and inference. Groups vary the training motion shift; underlining marks the best result within each group.}
% \label{tab:ablation_schedule}
% \begin{tabular}{cccc}
% \toprule
% Training $s_m$ & Inference $s_m$ & Jerk $\downarrow$ & FSR $\downarrow$\\
% \midrule
% \multirow{3}{*}{1} & 1 & \underline{1.251} & \underline{0.457}\\
%  & 10 & 8.003 & 0.986\\
%  & 15 & 9.703 & 1.000\\
% \midrule
% \multirow{3}{*}{10} & 1 & 7.021 & 0.985\\
%  & 10 & \underline{1.635} & \underline{0.635}\\
%  & 15 & 2.557 & 0.830\\
% \midrule
% \multirow{3}{*}{15} & 1 & 8.413 & 0.981\\
%  & 10 & 2.127 & 0.700\\
%  & 15 & \underline{1.809} & \underline{0.649}\\
% \bottomrule
% \end{tabular}
% \end{table}
% \FloatBarrier
% 

% Allow Figure 3 to share a page with the following text.
\section{Experiments}
\label{sec:experiments}
\subsection{Experimental Setup}
\label{sec:experimental-setup}

% Before metric-protocol clarification: \noindent\textbf{Dataset.} Our training dataset combines synthetic video--motion pairs from OMOMO~\citep{li2023omomo}, CHAIRs~\citep{jiang2023chairs}, and SAMP~\citep{hassan2021samp} with internet videos paired with CameraHMR~\citep{patel2025camerahmr} motion estimates. Section~\ref{sec:training-data} describes their construction. The synthetic component supplies known 3D motion, while the real component expands the range of environments and appearances. Text descriptions are generated using Gemini 2.5~\citep{comanici2025gemini}. We evaluate on a test set assembled from OMOMO, SAMP, COUCH~\citep{zhang2022couch}, and internet videos.
% Before dataset-and-splits revision: \noindent\textbf{Dataset.} Our training dataset combines synthetic video--motion pairs from OMOMO~\citep{li2023omomo}, CHAIRs~\citep{jiang2023chairs}, and SAMP~\citep{hassan2021samp} with internet videos paired with CameraHMR~\citep{patel2025camerahmr} motion estimates. Section~\ref{sec:training-data} describes their construction. The synthetic component supplies known 3D motion, while the real component expands the range of environments and appearances. Text descriptions are generated using Gemini 2.5~\citep{comanici2025gemini}. The main comparison uses 100 interaction cases: 12 from CHAIRs, 70 from COUCH~\citep{zhang2022couch}, 12 from OMOMO, and 6 from SAMP. The ablations use a separate 35-case set, described in Section~\ref{sec:ablation}.
% Before experiment citation revision: \noindent\textbf{Dataset Construction and Splits.} Following Section~\ref{sec:training-data}, HOI-mix combines SAMP, CHAIRs, and OMOMO motions paired with synthesized RGB videos and curated Bilibili videos with CameraHMR motion estimates. The training set contains 20,812 video--motion--text triplets totaling 29.15 hours, each with 121 frames at 24 fps. These triplets represent 10,847 distinct motion clips, with two appearance variants per interaction motion and one video per Bilibili clip. Motions share a common camera coordinate system; training uses sample-level shuffling across the merged dataset without source-specific reweighting. The legacy validation split contains 30 interaction and 20 non-interaction clips. Interaction data are held out by source sequence, and Bilibili data by clip. The final interaction benchmark contains 100 cases: 12 from CHAIRs, 12 from OMOMO, 6 from SAMP, and 70 additional cases from COUCH~\citep{zhang2022couch}. All benchmark source sequences are disjoint from the current training manifest. However, the benchmark retains the 30 interaction cases from the legacy validation split and is therefore not independent of that split. The separate 35-case ablation set is described in Section~\ref{sec:ablation}.
\noindent\textbf{Dataset Construction and Splits.} Following Section~\ref{sec:training-data}, HOI-mix combines SAMP~\citep{hassan2021samp}, CHAIRs~\citep{jiang2023chairs}, and OMOMO~\citep{li2023omomo} motions paired with synthesized RGB videos and curated Bilibili videos with CameraHMR~\citep{patel2025camerahmr} motion estimates. The training set contains 20,812 video--motion--text triplets totaling 29.15 hours, each with 121 frames at 24 fps. These triplets represent 10,847 distinct motion clips, with two appearance variants per interaction motion and one video per Bilibili clip. Motions share a common camera coordinate system; training uses sample-level shuffling across the merged dataset without source-specific reweighting. The legacy validation split contains 30 interaction and 20 non-interaction clips. Interaction data are held out by source sequence, and Bilibili data by clip. The final interaction benchmark contains 100 cases: 12 from CHAIRs, 12 from OMOMO, 6 from SAMP, and 70 additional cases from COUCH~\citep{zhang2022couch}. All benchmark source sequences are disjoint from the current training manifest. However, the benchmark retains the 30 interaction cases from the legacy validation split and is therefore not independent of that split. The separate 35-case ablation set is described in Section~\ref{sec:ablation}.

% Before experiment citation revision: \noindent\textbf{Implementation Details.} We initialize World2Motion from Cosmos 3 NANO (16B)~\citep{agarwal2026cosmos}, which includes a Qwen3-VL (8B) text backbone~\citep{bai2025qwen3}, 36 Mixture-of-Transformers layers, and a Wan2.2 causal VAE~\citep{wan2025wan}. The motion encoder and decoder are trained from scratch, and the understanding branch remains frozen. We train for 20,000 iterations with AdamW, a learning rate of $5\times10^{-5}$, a cosine decay schedule, and 200 warmup steps. Text conditioning is dropped with probability 0.1 for classifier-free guidance. Training uses bfloat16 precision, a batch size of 16, and 16 NVIDIA H100 GPUs, and takes approximately 44 hours.
\noindent\textbf{Implementation Details.} We initialize World2Motion from Cosmos 3 NANO (16B)~\citep{agarwal2026cosmos}, which includes a Qwen3-VL (8B) text backbone~\citep{bai2025qwen3}, 36 Mixture-of-Transformers layers, and a Wan2.2 causal VAE~\citep{wan2025wan}. The motion encoder and decoder are trained from scratch, and the understanding branch remains frozen. We train for 20,000 iterations with AdamW~\citep{loshchilov2019decoupled}, a learning rate of $5\times10^{-5}$, a cosine decay schedule, and 200 warmup steps. Text conditioning is dropped with probability 0.1 for classifier-free guidance~\citep{ho2022classifierfree}. Training uses bfloat16 precision, a batch size of 16, and 16 NVIDIA H100 GPUs, and takes approximately 44 hours.

% Before dataset-and-splits revision: \noindent\textbf{Baselines.} We compare with text-to-motion methods HY-Motion~\citep{hymotion2025}, Kimodo~\citep{Kimodo2026}, and Go-to-Zero~\citep{fan2025gotozero}; human--object interaction methods LIGHT~\citep{wang2026light} and HOIFHLI~\citep{wu2025hoifhli}; and a two-stage MiniMax-H3~\citep{minimaxh3} + CameraHMR pipeline. We also fine-tune CoMoVi~\citep{zhao2026comovi} on our training data as a joint video--motion baseline. Methods that accept an initial human pose---Kimodo, HOIFHLI, CoMoVi, and World2Motion---use the pose estimated by CameraHMR. This shared pose estimate can contain spatial errors that affect the subsequent motion. LIGHT and HOIFHLI additionally require explicit 3D object information; World2Motion obtains scene context from the input image.
% Before experiment citation revision: \noindent\textbf{Baselines.} We compare with text-to-motion methods HY-Motion~\citep{hymotion2025}, Kimodo~\citep{Kimodo2026}, and Go-to-Zero~\citep{fan2025gotozero}; interaction methods LIGHT~\citep{wang2026light} and HOIFHLI~\citep{wu2025hoifhli}; MiniMax-H3~\citep{minimaxh3} + CameraHMR; and CoMoVi~\citep{zhao2026comovi}, fine-tuned on our data. Kimodo, HOIFHLI, CoMoVi, and World2Motion use CameraHMR initial poses. LIGHT and HOIFHLI additionally require 3D object information; World2Motion uses the scene image.
\noindent\textbf{Baselines.} We compare with text-to-motion methods HY-Motion~\citep{hymotion2025}, Kimodo~\citep{Kimodo2026}, and Go-to-Zero~\citep{fan2025gotozero}; interaction methods LIGHT~\citep{wang2026light} and HOIFHLI~\citep{wu2025hoifhli}; MiniMax-H3~\citep{minimaxh3} + CameraHMR; and CoMoVi~\citep{zhao2026comovi}, fine-tuned on our data. We initialize Kimodo, HOIFHLI, CoMoVi, and World2Motion with poses estimated by CameraHMR~\citep{patel2025camerahmr}. LIGHT and HOIFHLI additionally require 3D object information, while World2Motion uses the scene image.

\subsection{Evaluation}
\label{sec:evaluation}

% Before metric-protocol clarification: \noindent\textbf{Motion Quality and Text Alignment.} We follow the MoMask~\citep{guo2024momask} evaluation protocol. Fr\'echet Inception Distance (FID) compares generated and reference motion feature distributions. R-Precision at rank 1 (R@1) and multimodal distance (MM-Dist) measure motion--text alignment. Lower FID and MM-Dist and higher R@1 are better. We also report jerk and foot sliding rate (FSR) to measure temporal irregularity and sliding artifacts, following WHAM~\citep{shin2024wham} and OmniControl~\citep{xie2024omnicontrol}; lower values are better for both.
% Before dataset-and-splits revision: \noindent\textbf{Motion Quality and Text Alignment.} We follow the MoMask~\citep{guo2024momask} evaluation protocol. Fr\'echet Inception Distance (FID) compares generated and reference motion feature distributions. R-Precision at rank 1 (R@1) and multimodal distance (MM-Dist) measure motion--text alignment. Lower FID and MM-Dist and higher R@1 are better. We report jerk and foot sliding rate (FSR), following WHAM~\citep{shin2024wham} and OmniControl~\citep{xie2024omnicontrol}. In Table~\ref{tab:comparison}, jerk uses each output's native frame rate and is reported in $10\,\mathrm{m/s^3}$. FSR is the percentage of adjacent frame pairs with sliding on either contacting foot, after resampling to 20 fps. Lower values are better. The ablations use a separate protocol, specified below.
\noindent\textbf{Motion Quality and Text Alignment.} Following MoMask~\citep{guo2024momask}, we report Fr\'echet Inception Distance (FID) for motion distributions, and R-Precision at rank 1 (R@1) and multimodal distance (MM-Dist) for motion--text alignment. Jerk and foot sliding rate (FSR) follow WHAM~\citep{shin2024wham} and OmniControl~\citep{xie2024omnicontrol}. In Table~\ref{tab:comparison}, jerk uses native frame rates and units of $10\,\mathrm{m/s^3}$; FSR measures the percentage of adjacent frame pairs with sliding on either contacting foot after resampling to 20 fps. Ablation protocols differ (Section~\ref{sec:ablation}); table arrows indicate the preferred direction.

% Before citation completion: \noindent\textbf{Scene Interaction.} Following OMOMO, we report Contact-F1. We render generated motions in the corresponding 3D scenes and use Qwen3-VL to assess whether the requested interaction is completed. Interaction Success Rate (ISR) is the fraction of test cases judged successful. These measures complement motion quality: a smooth motion can match the action description while failing to reach the intended object.
\noindent\textbf{Scene Interaction.} Following OMOMO~\citep{li2023omomo}, we report Contact-F1. We render generated motions in the corresponding 3D scenes and use Qwen3-VL~\citep{bai2025qwen3} to assess whether the requested interaction is completed. Interaction Success Rate (ISR) is the fraction of test cases judged successful. These measures complement motion quality: a smooth motion can match the action description while failing to reach the intended object.

% Before citation completion: \noindent\textbf{Video--Motion Consistency.} For joint generation methods, we recover motion from the generated video using CameraHMR and report PoseError against the directly generated motion. This measures agreement between the two outputs through a motion-recovery model, rather than error against a ground-truth 3D trajectory. We therefore interpret PoseError together with text alignment and scene interaction scores.
% Before metric-protocol clarification: \noindent\textbf{Video--Motion Consistency.} For joint generation methods, we recover motion from the generated video using CameraHMR~\citep{patel2025camerahmr} and report PoseError against the directly generated motion. This measures agreement between the two outputs through a motion-recovery model, rather than error against a ground-truth 3D trajectory. We therefore interpret PoseError together with text alignment and scene interaction scores.
\noindent\textbf{Video--Motion Consistency.} For joint generation methods, we recover motion from the generated video using CameraHMR~\citep{patel2025camerahmr} and compare it with the directly generated motion. PoseError is the pelvis-relative mean Euclidean distance over 22 corresponding body joints, in millimeters. We exclude invalid CameraHMR frames and average within each sequence, then across sequences. This measures cross-modal pose agreement, rather than ground-truth 3D accuracy or agreement of the root trajectories.

% Before dataset-and-splits revision: \noindent\textbf{Inference Efficiency.} We report mean inference time in seconds. World2Motion results in the main comparison use the post-processing described in Section~\ref{sec:method}. All ablation results are obtained without the Savitzky--Golay filter so that the effects of the motion representation and noise schedule can be assessed directly.
% Before experiment citation revision: \noindent\textbf{Inference Efficiency.} We report mean inference time in seconds. Main-comparison World2Motion outputs use Savitzky--Golay post-processing (Section~\ref{sec:method}); all ablations disable this filter.
\noindent\textbf{Inference Efficiency.} We report mean inference time in seconds. Main-comparison World2Motion outputs use Savitzky--Golay~\citep{savitzky1964smoothing} post-processing (Section~\ref{sec:method}); all ablations disable this filter.

% Previous quantitative preamble preserved; metrics are defined in Evaluation and the table caption.
% Table~\ref{tab:comparison} compares motion quality, scene interaction, video--motion consistency, and inference time across generation pipelines. We examine these aspects separately because the methods differ in their input requirements and whether they produce both video and motion.
% Continue after the setup without creating a mostly empty page.
\subsection{Quantitative Comparisons}
\begin{table}[!htbp]
\centering
\small
% Before metric-protocol clarification: \caption{Quantitative comparison on the multi-source test set. FID measures motion distribution quality; R@1 and MM-Dist measure text alignment; Contact-F1 and ISR evaluate scene interaction. Jerk and FSR measure motion stability, and PoseError measures video--motion agreement. Time is mean inference time in seconds. Bold denotes the best reported value per column; dashes indicate unreported scores.}
\caption{Quantitative comparison on 100 interaction cases. FID measures motion distribution quality; R@1 and MM-Dist measure text alignment; Contact-F1 and ISR evaluate scene interaction. Jerk and FSR measure motion stability, and PoseError measures video--motion agreement. Time is mean inference time in seconds. Bold denotes the best reported value per column; dashes indicate unreported scores.}
\label{tab:comparison}
\setlength{\tabcolsep}{2pt}
\fitwidth{\textwidth}{%
\begin{tabular}{l c c c c c c c c c}
\toprule
% Before metric-protocol clarification: Method & FID$\downarrow$ & R@1$\uparrow$ & MM-Dist$\downarrow$ & Contact-F1$\uparrow$ & ISR$\uparrow$ & Jerk$\downarrow$ & FSR$\downarrow$ & PoseError$\downarrow$ & Time (s)$\downarrow$\\
Method & FID$\downarrow$ & R@1$\uparrow$ & MM-Dist$\downarrow$ & Contact-F1$\uparrow$ & ISR$\uparrow$ & \begin{tabular}[c]{@{}c@{}}Jerk$\downarrow$\\($10\,\mathrm{m/s^3}$)\end{tabular} & \begin{tabular}[c]{@{}c@{}}FSR$\downarrow$\\(\%)\end{tabular} & \begin{tabular}[c]{@{}c@{}}PoseError$\downarrow$\\(mm)\end{tabular} & Time (s)$\downarrow$\\
\midrule
\multicolumn{10}{l}{\textbf{Text-to-motion}}\\
% Before citation completion: HY-Motion & 0.46 & 0.10 & 1.21 & 0.09 & 0.51 & 3.43 & \textbf{1.25} & -- & 2.83\\
HY-Motion~\citep{hymotion2025} & 0.46 & 0.10 & 1.21 & 0.09 & 0.51 & 3.43 & \textbf{1.25} & -- & 2.83\\
% Before citation completion: Kimodo & 0.66 & 0.05 & 1.30 & 0.03 & 0.23 & 35.14 & 6.07 & -- & \textbf{1.93}\\
Kimodo~\citep{Kimodo2026} & 0.66 & 0.05 & 1.30 & 0.03 & 0.23 & 35.14 & 6.07 & -- & \textbf{1.93}\\
% Before citation completion: Go-to-Zero & 0.51 & 0.04 & 1.35 & 0.07 & 0.36 & 17.50 & 3.42 & -- & 2.57\\
Go-to-Zero~\citep{fan2025gotozero} & 0.51 & 0.04 & 1.35 & 0.07 & 0.36 & 17.50 & 3.42 & -- & 2.57\\
\midrule
\multicolumn{10}{l}{\textbf{Human--object interaction}}\\
% Before citation completion: HOIFHLI & 0.82 & 0.02 & 1.37 & 0.09 & 0.12 & 10.62 & 6.28 & -- & 13.94\\
HOIFHLI~\citep{wu2025hoifhli} & 0.82 & 0.02 & 1.37 & 0.09 & 0.12 & 10.62 & 6.28 & -- & 13.94\\
% Before citation completion: LIGHT & 0.62 & 0.03 & 1.35 & 0.17 & 0.46 & \textbf{1.37} & 11.60 & -- & 25.88\\
LIGHT~\citep{wang2026light} & 0.62 & 0.03 & 1.35 & 0.17 & 0.46 & \textbf{1.37} & 11.60 & -- & 25.88\\
\midrule
\multicolumn{10}{l}{\textbf{Two-stage video generation and motion recovery}}\\
% Before citation completion: \begin{tabular}[c]{@{}l@{}}MiniMax-H3 +\\CameraHMR\end{tabular} & 0.69 & 0.10 & 1.32 & \textbf{0.24} & \textbf{0.91} & 37.22 & 6.41 & -- & 523.67\\
\begin{tabular}[c]{@{}l@{}}MiniMax-H3~\citep{minimaxh3} +\\CameraHMR~\citep{patel2025camerahmr}\end{tabular} & 0.69 & 0.10 & 1.32 & \textbf{0.24} & \textbf{0.91} & 37.22 & 6.41 & -- & 523.67\\
\midrule
\multicolumn{10}{l}{\textbf{Joint video--motion generation}}\\
% Before citation completion: CoMoVi & 0.71 & 0.12 & 1.25 & 0.09 & 0.69 & 16.10 & 24.45 & 81.28 & 270.45\\
CoMoVi~\citep{zhao2026comovi} & 0.71 & 0.12 & 1.25 & 0.09 & 0.69 & 16.10 & 24.45 & 81.28 & 270.45\\
World2Motion & \textbf{0.27} & \textbf{0.24} & \textbf{1.13} & 0.22 & \textbf{0.91} & 8.35 & 2.80 & \textbf{67.50} & 156.59\\
\bottomrule
\end{tabular}
}
\end{table}

% Before concise result discussion: \noindent\textbf{Motion Quality and Text Alignment.} World2Motion achieves the lowest FID and MM-Dist and the highest R@1 among the compared methods. Relative to HY-Motion, FID decreases from 0.46 to 0.27 and R@1 increases from 0.10 to 0.24. The improvement is also observed against CoMoVi, which produces both video and motion but has an FID of 0.71 and R@1 of 0.12. These results show improved alignment between the generated 3D motion and the action prompt, alongside the generation of a corresponding video.
% Before experiment citation revision: \noindent\textbf{Motion Quality and Text Alignment.} World2Motion achieves the lowest FID and MM-Dist and the highest R@1 among the compared methods. Relative to HY-Motion, FID decreases from 0.46 to 0.27 and R@1 increases from 0.10 to 0.24. It also improves over CoMoVi, whose FID and R@1 are 0.71 and 0.12, respectively.
\noindent\textbf{Motion Quality and Text Alignment.} In Table~\ref{tab:comparison}, World2Motion achieves the lowest FID and MM-Dist and the highest R@1 among the compared methods. Relative to HY-Motion, FID decreases from 0.46 to 0.27 and R@1 increases from 0.10 to 0.24. It also improves over CoMoVi, whose FID and R@1 are 0.71 and 0.12, respectively.

\noindent\textbf{Scene Interaction.} World2Motion reaches an ISR of 0.91, compared with 0.23--0.51 for the evaluated text-to-motion methods and 0.12--0.46 for the human--object interaction methods. Its Contact-F1 of 0.22 also exceeds those of these baselines. The two-stage pipeline achieves the same ISR and a higher Contact-F1 of 0.24. Successful completion and contact accuracy are therefore distinct: matching the overall interaction success of the two-stage method does not imply identical contact quality. The remaining contact gap motivates more precise supervision for human--object interactions.

\noindent\textbf{Motion Stability.} World2Motion has lower jerk and FSR than Kimodo, Go-to-Zero, HOIFHLI, and the two video-based baselines. HY-Motion achieves lower values on both stability metrics, while LIGHT achieves the lowest jerk but has a higher FSR of 11.60. Thus, the best result depends on which aspect of stability is measured. World2Motion improves stability over the compared video-based pipelines, but it does not surpass every motion-only method. We retain both jerk and FSR to distinguish temporal irregularity from sliding artifacts.

\noindent\textbf{Video--Motion Consistency.} Compared with CoMoVi, World2Motion reduces PoseError from 81.28 to 67.50, a 17.0\% reduction. This indicates closer agreement between its generated video and directly predicted 3D motion under the motion-recovery evaluation. It complements the improvements in text alignment and scene interaction, which assess the generated motion against the prompt and environment rather than against the other output modality.

\noindent\textbf{Inference Efficiency.} World2Motion takes 156.59 seconds on average, compared with 523.67 seconds for MiniMax-H3 + CameraHMR and 270.45 seconds for CoMoVi: speedups of approximately $3.34\times$ and $1.73\times$, respectively. The text-to-motion and human--object interaction methods remain faster. Our efficiency gain concerns video-based motion generation. World2Motion provides both a video and an explicit 3D motion sequence, whereas the faster motion-only baselines do not produce the corresponding video. The results support reducing the cost of obtaining these paired outputs while preserving competitive interaction performance.

% Continuous preprint layout: no forced page break.
\subsection{Qualitative Comparisons}

\begin{figure}[!htbp]
\centering
\includegraphics[width=0.96\linewidth]{Fig/pdf/Comparison_v2.pdf}
\caption{Qualitative comparison of scene interactions: (a) lifting a plastic box, (b) walking to a chair and sitting, and (c) walking to a sofa and sitting. Each method shows successive motion states; video-generating methods also show corresponding video frames. Human--object interaction baselines include their predicted object trajectories.}
\label{fig:Comparison}
\end{figure}

Figure~\ref{fig:Comparison} compares actions that depend on both body motion and placement relative to an object. In the sitting examples, producing a plausible sitting pose is insufficient: the character must first approach the intended seat and then align the body with it. The box example also involves coordinating the body with object motion. Viewing the approach and final interaction together helps distinguish action matching from scene-consistent execution. For video-based methods, the paired rows additionally allow the reader to inspect whether the generated video and recovered or directly generated 3D motion describe the same interaction.

% Continuous preprint layout: no forced page break.
\subsection{Ablation Studies}
\label{sec:ablation}

% Before metric-protocol clarification: We evaluate root translation and noise scheduling using jerk and FSR. All ablations disable post-processing. Comparisons within each study keep other settings fixed. The first study compares motion representations; the second changes training and inference shifts independently. This separates the choice of noise schedule from the effect of using it consistently.
% Before condensing discussion to accommodate Limitations: We evaluate root translation and noise scheduling using 10,000-iteration checkpoints on a separate 35-case set (15 demonstration clips and 20 internet-video clips), without post-processing. Table~\ref{tab:comparison} instead uses the 20,000-iteration model with smoothing on 100 interaction cases. Ablation jerk uses 25 fps and is reported in $10^3\,\mathrm{m/s^3}$. Ablation FSR is the fraction of contacting foot-joint/frame pairs that slide, pooled over four foot joints on a normalized skeleton, in the range 0--1. Its denominator differs from the all-frame-pair denominator in Table~\ref{tab:comparison}; the scores are therefore compared only within the ablation protocol. The relative-translation row in Table~\ref{tab:ablation_norm} and the matched $s_m=1$ entry in Table~\ref{tab:ablation_schedule} share the same generated outputs.
All ablations use 10,000-iteration checkpoints and 35 clips (15 demonstration clips and 20 internet-video clips), without post-processing; Table~\ref{tab:comparison} uses the smoothed 20,000-iteration model on 100 interaction cases. Ablation jerk uses 25 fps. FSR pools sliding observations over contacting foot-joint/frame pairs from four foot joints on a normalized skeleton. Unlike Table~\ref{tab:comparison}, its denominator excludes non-contact observations, so the scores are not directly comparable across protocols. The relative-translation row in Table~\ref{tab:ablation_norm} and the matched $s_m=1$ entry in Table~\ref{tab:ablation_schedule} share the same outputs.

% Previous root-translation discussion and table; full table moved to Appendix.
% \noindent\textbf{Root translation.} Representing each root position relative to the initial position reduces jerk from 4.355 to 1.251 (71.3\%) and FSR from 0.766 to 0.457 (40.3\%), compared with absolute translation (Table~\ref{tab:ablation_norm}). Both representations retain root position and 6D joint rotations; the change is the reference position used for translation. The comparison therefore isolates this coordinate choice while keeping the rotation representation and other settings fixed. These results support using relative positions for more stable generation without requiring root velocity as an input.
% 
% \begin{table}[!htbp]
% \centering
% \small
% \setlength{\tabcolsep}{5pt}
% \caption{Root translation ablation without post-processing. Relative positions are measured from the initial root position. Bold marks the lower value in each column.}
% \label{tab:ablation_norm}
% \begin{tabular}{lcc}
% \toprule
% Root translation & Jerk $\downarrow$ & FSR $\downarrow$\\
% \midrule
% Absolute & 4.355 & 0.766\\
% Relative to initial position & \textbf{1.251} & \textbf{0.457}\\
% \bottomrule
% \end{tabular}
% \end{table}
% Previous condensed root-translation discussion preserved.
% \noindent\textbf{Root translation.} With other settings fixed and post-processing disabled, expressing root positions relative to the initial position reduces jerk from 4.355 to 1.251 (71.3\%) and FSR from 0.766 to 0.457 (40.3\%), compared with absolute positions (Appendix Table~\ref{tab:ablation_norm}).
\begin{wraptable}{r}{0.45\textwidth}
\centering
\small
% Before metric-protocol clarification: \caption{Root translation ablation without post-processing. Relative positions use the initial root position as the reference.}
% Before qualified ablation metric labels: \caption{Root translation ablation on 35 cases, without post-processing. Relative positions use the initial root position as the reference.}
\setlength{\belowcaptionskip}{6pt}
\caption{Root translation ablation on 35 cases without post-processing. Raw Jerk uses 25 fps ($10^3\,\mathrm{m/s^3}$). Contact FSR is the sliding fraction among contacting foot-joint/frame pairs on a normalized skeleton (0--1); its definition differs from Table~\ref{tab:comparison}.}
\label{tab:ablation_norm}
\setlength{\tabcolsep}{4pt}
\begin{tabular}{l c c}
\toprule
% Before metric-protocol clarification: Translation & Jerk $\downarrow$ & FSR $\downarrow$\\
% Before qualified ablation metric labels: Translation & \begin{tabular}[c]{@{}c@{}}Jerk $\downarrow$\\($10^3\,\mathrm{m/s^3}$)\end{tabular} & \begin{tabular}[c]{@{}c@{}}FSR $\downarrow$\\(0--1)\end{tabular}\\
Translation & \begin{tabular}[c]{@{}c@{}}Raw Jerk $\downarrow$\\($10^3\,\mathrm{m/s^3}$)\end{tabular} & \begin{tabular}[c]{@{}c@{}}Contact FSR $\downarrow$\\(0--1)\end{tabular}\\
\midrule
Absolute & 4.355 & 0.766\\
Relative & \textbf{1.251} & \textbf{0.457}\\
\bottomrule
\end{tabular}
\end{wraptable}
% Before condensing discussion to accommodate Limitations: \noindent\textbf{Root translation.} Table~\ref{tab:ablation_norm} compares absolute root positions with positions relative to the initial frame. With other settings fixed and post-processing disabled, relative translation reduces jerk from 4.355 to 1.251 (71.3\%) and FSR from 0.766 to 0.457 (40.3\%). Both configurations use the same 6D joint rotations; only the translation reference changes. The improvements therefore support this coordinate choice for temporal stability, without introducing root velocity features or relying on smoothing.
\noindent\textbf{Root translation.} With other settings fixed, relative root translation reduces jerk from 4.355 to 1.251 (71.3\%) and FSR from 0.766 to 0.457 (40.3\%) compared with absolute translation (Table~\ref{tab:ablation_norm}). Both configurations use the same 6D joint-rotation representation. These results support the relative coordinate choice for temporal stability in the tested setting, without root velocity features or post-processing.

% Before condensing discussion to accommodate Limitations: \noindent\textbf{Shift-decoupled noise schedule.} Table~\ref{tab:ablation_schedule} separates two questions: which shifts to use, and whether the training and inference shifts should match. With matched shifts, $s_m=1$ achieves lower jerk and FSR than $s_m=10$ or $15$. Compared with matched $s_m=10$, matched $s_m=1$ reduces jerk by 23.5\% and FSR by 28.0\%. For each training configuration, matching the inference shift also gives the lowest values on both metrics. For example, after training with $s_m=1$, changing only the inference shift to $10$ increases jerk from 1.251 to 8.003. These observations support modality-specific shifts and consistent training and inference noise relationships. They do not require an additional smoothing step, since the filter is disabled throughout this comparison.
\noindent\textbf{Shift-decoupled noise schedule.} In Table~\ref{tab:ablation_schedule}, matched training and inference shifts of $s_m=1$ reduce jerk and FSR by 23.5\% and 28.0\% relative to matched $s_m=10$. Matching shifts gives the lowest values within every row. After training with $s_m=1$, changing only the inference shift to $10$ increases jerk from 1.251 to 8.003. These results support modality-specific shifts and training--inference consistency for the tested stability metrics.

% Previous nine-row schedule table preserved.
% \begin{table}[!htbp]
% \centering
% \small
% \setlength{\tabcolsep}{5pt}
% \caption{Noise-schedule ablation without post-processing. Video shift $s_v=10$ is fixed in training and inference. Groups vary the training motion shift; bold marks the best result within each group.}
% \label{tab:ablation_schedule}
% \begin{tabular}{cccc}
% \toprule
% Training $s_m$ & Inference $s_m$ & Jerk $\downarrow$ & FSR $\downarrow$\\
% \midrule
% \multirow{3}{*}{1} & 1 & \textbf{1.251} & \textbf{0.457}\\
%  & 10 & 8.003 & 0.986\\
%  & 15 & 9.703 & 1.000\\
% \midrule
% \multirow{3}{*}{10} & 1 & 7.021 & 0.985\\
%  & 10 & \textbf{1.635} & \textbf{0.635}\\
%  & 15 & 2.557 & 0.830\\
% \midrule
% \multirow{3}{*}{15} & 1 & 8.413 & 0.981\\
%  & 10 & 2.127 & 0.700\\
%  & 15 & \textbf{1.809} & \textbf{0.649}\\
% \bottomrule
% \end{tabular}
% \end{table}
% Previous two-block schedule layout preserved.
% \begin{table}[!htbp]
% \centering
% \small
% \caption{Noise-schedule ablation without post-processing. Video shift $s_v=10$ is fixed in training and inference. Bold marks the lowest value in each training row.}
% \label{tab:ablation_schedule}
% \setlength{\tabcolsep}{5pt}
% \begin{minipage}[t]{0.48\linewidth}
% \centering
% \textbf{Jerk $\downarrow$}\par\smallskip
% \begin{tabular}{c c c c}
% \toprule
%  & \multicolumn{3}{c}{Inference $s_m$}\\
% \cmidrule(lr){2-4}
% Training $s_m$ & 1 & 10 & 15\\
% \midrule
% 1 & \textbf{1.251} & 8.003 & 9.703\\
% 10 & 7.021 & \textbf{1.635} & 2.557\\
% 15 & 8.413 & 2.127 & \textbf{1.809}\\
% \bottomrule
% \end{tabular}
% \end{minipage}\hfill
% \begin{minipage}[t]{0.48\linewidth}
% \centering
% \textbf{FSR $\downarrow$}\par\smallskip
% \begin{tabular}{c c c c}
% \toprule
%  & \multicolumn{3}{c}{Inference $s_m$}\\
% \cmidrule(lr){2-4}
% Training $s_m$ & 1 & 10 & 15\\
% \midrule
% 1 & \textbf{0.457} & 0.986 & 1.000\\
% 10 & 0.985 & \textbf{0.635} & 0.830\\
% 15 & 0.981 & 0.700 & \textbf{0.649}\\
% \bottomrule
% \end{tabular}
% \end{minipage}
% \end{table}
\begin{table}[!htbp]
\centering
\small
% Before metric-protocol clarification: \caption{Noise-schedule ablation without post-processing. Columns within each metric group specify the inference motion shift $s_m$. Video shift $s_v=10$ is fixed throughout. Bold marks the lowest value for each metric within each row.}
% Before qualified ablation metric labels: \caption{Noise-schedule ablation on the same 35 cases as Table~\ref{tab:ablation_norm}, without post-processing. Columns within each metric group specify the inference motion shift $s_m$. Video shift $s_v=10$ is fixed throughout. Bold marks the lowest value for each metric within each row.}
\setlength{\belowcaptionskip}{6pt}
\caption{Noise-schedule ablation on the same 35 cases, using Raw Jerk and contact-conditioned FSR as defined in Table~\ref{tab:ablation_norm}. These unfiltered ablation scores are not directly comparable to Table~\ref{tab:comparison}. Columns give the inference motion shift $s_m$; $s_v=10$ is fixed. Bold marks the lowest value per metric within each row.}
\label{tab:ablation_schedule}
\setlength{\tabcolsep}{7pt}
\begin{tabular}{c c c c c c c}
\toprule
% Before metric-protocol clarification:  & \multicolumn{3}{c}{Jerk $\downarrow$} & \multicolumn{3}{c}{FSR $\downarrow$}\\
% Before qualified ablation metric labels:  & \multicolumn{3}{c}{Jerk ($10^3\,\mathrm{m/s^3}$) $\downarrow$} & \multicolumn{3}{c}{FSR (0--1) $\downarrow$}\\
 & \multicolumn{3}{c}{Raw Jerk ($10^3\,\mathrm{m/s^3}$) $\downarrow$} & \multicolumn{3}{c}{Contact FSR (0--1) $\downarrow$}\\
\cmidrule(lr){2-4}\cmidrule(lr){5-7}
Training $s_m$ & 1 & 10 & 15 & 1 & 10 & 15\\
\midrule
1 & \textbf{1.251} & 8.003 & 9.703 & \textbf{0.457} & 0.986 & 1.000\\
10 & 7.021 & \textbf{1.635} & 2.557 & 0.985 & \textbf{0.635} & 0.830\\
15 & 8.413 & 2.127 & \textbf{1.809} & 0.981 & 0.700 & \textbf{0.649}\\
\bottomrule
\end{tabular}
\end{table}
% Let floats share the following text page.
% Previous conclusion preserved (2026-09-26).
% \section{Conclusion}
% We present Tandem, a MoT-based framework that synchronously generates videos and motions. We design a coupled noise schedule and a motion representation to improve the motion quality and reduce the jitter. Our ablation studies validate the effectiveness of our strategies. The comparison results with other motion generation methods demonstrate the convenience, scene-awareness, speed and generalization of our framework. Our method benefits the game development by automatically generating character motions, reducing the labor-intensive Keyframe animations. However, our framework can't generate motions in real time and synthesize fine-grained hand interactions, which limits the further application in the game. We will curate hand-interaction datasets and use distillation to achieve faster and detailed interactions in the future.
% \label{sec:conclusion}
% 
% 


% Previous conclusion preserved before separate Limitations section.
% \section{Conclusion}
% \label{sec:conclusion}
% We present World2Motion, which adapts Cosmos 3 to jointly generate scene-aware 3D human motion and video from an image and a text prompt. A hybrid dataset supports this adaptation, while relative root translation and a shift-decoupled noise schedule improve temporal stability. Experiments show improved motion--text alignment and competitive scene interaction, with approximately $3.34\times$ faster inference than the evaluated two-stage pipeline. Contact accuracy remains a limitation, as the two-stage baseline retains a higher Contact-F1. Generation remains slower than real time, and fine-grained hand interactions remain challenging. Future work includes improving inference efficiency and providing more precise supervision for hand--object contact.

\section{Limitations}
\label{sec:limitations}
% Before clarifying data coverage and evaluation limits: We view World2Motion as an initial baseline for adapting pretrained video world models to 3D human motion generation. Our training data cover a limited range of interactions, and performance on more complex cases has not been systematically evaluated. CameraHMR errors can affect both initial poses and supervision from real videos. Inference is faster than CoMoVi in our comparison, but remains too slow for real-time applications. Reliable root placement and precise contact remain challenges, with Contact-F1 below the two-stage baseline. Richer interaction data, additional modalities such as scene geometry and object states, and faster generation are directions for future work. For current offline applications, generated motions can initialize agent-assisted workflows for inspection and refinement. The supplementary video illustrates such a workflow, including motion smoothing and contact refinement. These examples illustrate a use of our outputs; broader gains in workflow quality and efficiency require systematic evaluation.
% Previous Limitations opening: We view World2Motion as an initial baseline for adapting pretrained video world models to 3D human motion generation.
% Before author-approved rewrite: World2Motion provides a baseline for adapting pretrained video world models to 3D human motion generation. Training data mainly cover single-person motion and simple interactions; more complex cases lack systematic evaluation. CameraHMR errors can affect initial poses and real-video supervision. Inference is faster than CoMoVi in our comparison, but remains unsuitable for real-time applications. Our metrics provide a limited assessment: ISR and PoseError do not guarantee precise contact. Root placement and contact remain challenges, with Contact-F1 below the two-stage baseline. Future directions include richer interactions, additional modalities such as scene geometry and object states, and faster generation. For offline applications, our outputs can initialize agent-assisted inspection and refinement. The supplementary video illustrates motion smoothing and contact refinement in such a workflow; broader quality and efficiency gains require systematic evaluation.
World2Motion provides a baseline for adapting pretrained video world models to 3D human motion generation, but several limitations remain. First, our training data mainly cover single-person motion and simple interactions, leaving more complex scenarios largely unexplored. Second, despite the improvement in inference speed, World2Motion is still too slow for real-time use. Third, accurate root placement and contact remain challenging. Our evaluation also has limits: a high ISR or low PoseError does not imply precise contact. Future work could expand interaction coverage, incorporate scene geometry and object states, and reduce inference time. For now, generated motions can serve as starting points for further editing. The supplementary video shows examples of agent-assisted smoothing and contact refinement.

\section{Conclusion}
\label{sec:conclusion}
% Before removing the repeated baseline-positioning sentence: World2Motion adapts Cosmos 3 to jointly generate scene-aware 3D human motion and video from an image and a text prompt. Our adaptation combines hybrid video--motion training data, relative root translation, and a shift-decoupled noise schedule. The ablations show reduced jerk and foot sliding under the tested settings. This work provides a baseline for further study of video-model adaptation to image-conditioned motion generation.
% Before author-approved rewrite: World2Motion adapts Cosmos 3 to jointly generate scene-aware 3D human motion and video from an image and a text prompt. Our adaptation combines hybrid video--motion training data, relative root translation, and a shift-decoupled noise schedule. The ablations show reduced jerk and foot sliding under the tested settings.
We present World2Motion, which adapts Cosmos 3 to jointly generate scene-aware 3D human motion and video from a single image and a text prompt. We train the model on synthetic video--motion pairs and real videos with estimated motion, and use a shift-decoupled noise schedule to reduce motion jitter. On our interaction benchmark, World2Motion improves motion--text alignment over the compared methods and matches the interaction success rate of the two-stage baseline.


% Continuous preprint layout: no forced page break.
% Unfilled statement templates preserved; omitted from the PDF.
% \subsection*{AI use statement}
% 
% In this work, we used generative AI tools for [tasks with required disclosure].
% We have not used generative AI tools for [other tasks with required disclosure],
% and [the rest of the required disclosure tasks] are not applicable to this work.
% Additionally, we used generative AI tools for [tasks with recommended
% disclosure]. We have reviewed all AI-assisted work. We take responsibility for
% the final content of this work, including text, claims or artifacts produced
% with the aid of generative AI.
% 
% \subsection*{Ethics statement}
% 
% TODO
% 
% \subsection*{Reproducibility statement}
% 
% TODO
% 
% 
% 
% \subsection*{AI use statement}
% Generative AI tools were used to assist with manuscript editing and organization. Uses of generative models in training-data construction and evaluation are described in Sections~\ref{sec:training-data} and~\ref{sec:experimental-setup}.

\FloatBarrier
\bibliography{iclr2027_conference}
\bibliographystyle{iclr2027_conference}

% Unwritten appendix outline preserved; omitted from the PDF.
% Continuous preprint layout: no forced page break.
\clearpage
\appendix
% Previous appendix heading: \section{Appendix}
\section{Additional Qualitative Results}

Figure~\ref{fig:supplementary-results-1} presents box lifting and carrying, approaching and sitting on a stool or armchair, and two yoga movements. The sitting examples show the approach to a seat as well as the transition from standing to sitting, allowing the reader to inspect body placement throughout the sequence. The yoga examples show full-body pose changes without object manipulation.


% Root-translation table moved back to the main text as a wraptable.
% \begin{table}[!htbp]
% \centering
% \small
% \setlength{\tabcolsep}{5pt}
% \caption{Root translation ablation without post-processing. Relative positions are measured from the initial root position. Bold marks the lower value in each column.}
% \label{tab:ablation_norm}
% \begin{tabular}{lcc}
% \toprule
% Root translation & Jerk $\downarrow$ & FSR $\downarrow$\\
% \midrule
% Absolute & 4.355 & 0.766\\
% Relative to initial position & \textbf{1.251} & \textbf{0.457}\\
% \bottomrule
% \end{tabular}
% \end{table}


\begin{figure}[!htbp]
    \centering
    % Previous figure width: \linewidth
    \includegraphics[width=0.90\linewidth]{Fig/pdf/Supplementary_Results_01_v3.pdf}
        % Previous caption: \caption{Overview of World2Motion. From an image and a text prompt, the adapted Cosmos 3 backbone jointly generates 3D human motion and video. CameraHMR estimates the initial pose from the image. The shift-decoupled schedule uses modality-specific noise shifts with shared denoising progress.}
    \caption{Qualitative results of World2Motion for box lifting and carrying, sitting on a stool or armchair, and two yoga movements. Each example shows five successive video frames (top) and corresponding 3D motion renderings (bottom), ordered from left to right. The input text prompt appears below each example.}
    % Previous label: \label{fig:Method}
    \label{fig:supplementary-results-1}
\end{figure}

\clearpage
Figure~\ref{fig:supplementary-results-2} adds sofa sitting, picking up and carrying a vacuum cleaner, standing up from a stool, and two dance sequences. These examples include both changes in body position relative to furniture and coordinated limb movements. In both figures, the paired video and motion rows show five corresponding time points, making it possible to compare how the requested action progresses in the two outputs. The text prompt is provided below each example.

\begin{figure}[!htbp]
    \centering
    % Previous figure width: \linewidth
    \includegraphics[width=0.90\linewidth]{Fig/pdf/Supplementary_Results_02_v3.pdf}
        % Previous caption: \caption{Overview of World2Motion. From an image and a text prompt, the adapted Cosmos 3 backbone jointly generates 3D human motion and video. CameraHMR estimates the initial pose from the image. The shift-decoupled schedule uses modality-specific noise shifts with shared denoising progress.}
    \caption{Additional qualitative results of World2Motion for sofa sitting, picking up and carrying a vacuum cleaner, standing up from a stool, and two dance sequences. Each example shows five successive video frames (top) and corresponding 3D motion renderings (bottom), ordered from left to right. The input text prompt appears below each example.}
    % Previous label: \label{fig:Method}
    \label{fig:supplementary-results-2}
\end{figure}



\end{document}
